\documentclass[11pt]{article}

\usepackage[utf8]{inputenc}
% microtype 의 폰트 확장은 scalable 폰트를 요구한다. 기본 CM 이 비트맵(Type3)으로
% 잡히는 배포판에서는 lmodern 없이는 "auto expansion is only possible with
% scalable fonts" 로 컴파일이 죽는다.
\usepackage{lmodern}
\usepackage[T1]{fontenc}
\usepackage[margin=1in]{geometry}
\usepackage{amsmath,amssymb}
\usepackage{graphicx}
\usepackage{longtable,booktabs}
\usepackage{microtype}
% svgnames 를 줘야 teal 이 정의된다. 없으면 \cite 하나마다
% "Undefined color `teal'" 이 뜬다 (본편 산출물에서 732회).
\usepackage[svgnames]{xcolor}
% hyperref goes last; it makes \cite clickable through to the bibliography.
\usepackage[colorlinks=true,linkcolor=blue,citecolor=teal,urlcolor=magenta]{hyperref}

% pandoc emits \tightlist but expects the preamble to define it.
\providecommand{\tightlist}{%
  \setlength{\itemsep}{0pt}\setlength{\parskip}{0pt}}

% If the compile times out on Overleaf, comment out \tableofcontents below first.
\title{Comprehensive Survey of Edge Computing: Architecture, Applications, and Challenges}
\author{Generated by AutoSurvey}
\date{}

\begin{document}
\maketitle
\tableofcontents
\newpage

\section{Introduction to Edge Computing}

\subsection{Definition and History of Edge Computing}

Edge computing is a novel distributed computing paradigm that extends cloud resources to the edge of the network, enabling data processing and analysis closer to where it is generated \cite{ref1}. This paradigm has emerged as a response to the growing need for reduced latency, improved real-time processing capabilities, and enhanced user experience. By processing data closer to the source, edge computing can significantly decrease the time it takes for data to travel from the device to the cloud and back, resulting in faster response times and improved overall performance \cite{ref2}.

The evolution of edge computing has been marked by several key milestones, including the introduction of cloudlet-based edge computing, which aimed to provide a more efficient and scalable way of delivering cloud services to edge devices \cite{ref2}. Cloudlets are small-scale data centers that are deployed at the edge of the network, providing a cache of cloud services and enabling faster access to applications and data. The concept of cloudlets has since evolved to include other edge computing paradigms, such as fog computing and mobile edge computing \cite{ref3}. Fog computing, introduced by Cisco in 2012, is an edge computing paradigm that extends cloud computing to the edge of the network, using a distributed architecture to provide low-latency and location-aware services \cite{ref1}. Mobile edge computing, on the other hand, is a paradigm that focuses on providing computing resources and services at the edge of mobile networks, enabling faster and more efficient processing of mobile data \cite{ref3}.

The current state of edge computing is characterized by a growing interest in the development of edge-based applications and services, driven by the increasing demand for low-latency and real-time processing capabilities \cite{ref4}. Edge computing is being applied in various domains, including industrial automation, smart cities, healthcare, and transportation, among others \cite{ref2}. The use of edge computing in these domains is enabling new use cases and applications, such as predictive maintenance, smart traffic management, and remote healthcare monitoring \cite{ref3}. As edge computing continues to evolve, it is likely to have a significant impact on various industries and applications, leading to increased productivity, safety, and profitability.

Despite the growing interest in edge computing, there are still several challenges and open research directions that need to be addressed, including scalability, reliability, and standardization \cite{ref1}. The development of edge computing standards and frameworks is crucial to enable interoperability and seamless communication between different edge devices and systems \cite{ref2}. Additionally, there is a need for more research on edge computing security, privacy, and trust, as well as on the development of edge-based machine learning and artificial intelligence applications \cite{ref4}. The integration of edge computing with artificial intelligence and machine learning is a particularly promising area of research, with potential applications in areas such as edge-based deep learning, reinforcement learning, and federated learning \cite{ref5}.

In recent years, there has been a growing interest in the development of edge-based machine learning and artificial intelligence applications, driven by the increasing demand for low-latency and real-time processing capabilities \cite{ref6}. Edge-based machine learning and artificial intelligence are enabling new use cases and applications, such as edge-based deep learning, reinforcement learning, and federated learning \cite{ref5}. However, there are still several challenges and open research directions that need to be addressed, including the development of edge-based machine learning frameworks and the optimization of edge-based machine learning algorithms \cite{ref7}. As edge computing continues to evolve, it is likely to play an increasingly important role in enabling faster and more reliable access to data and applications, and its benefits will be felt across various industries and applications \cite{ref1}.

In conclusion, edge computing is a novel distributed computing paradigm that extends cloud resources to the edge of the network, enabling data processing and analysis closer to where it is generated \cite{ref1}. The history of edge computing is characterized by the introduction of cloudlet-based edge computing, fog computing, and mobile edge computing, among other paradigms \cite{ref2}. The current state of edge computing is driven by the increasing demand for low-latency and real-time processing capabilities, and is being applied in various domains, including industrial automation, smart cities, healthcare, and transportation, among others \cite{ref3}. As edge computing continues to evolve, it is likely to have a significant impact on various industries and applications, leading to increased productivity, safety, and profitability, and its benefits will be discussed in more detail in the following section.

\subsection{Importance and Benefits of Edge Computing}

The importance and benefits of edge computing are multifaceted, and its impact is being felt across various industries and applications. As discussed in the previous section, edge computing has emerged as a response to the growing need for reduced latency, improved real-time processing capabilities, and enhanced user experience. One of the primary advantages of edge computing is its ability to reduce latency, which is critical for real-time processing and decision-making \cite{ref2}. By processing data closer to the source, edge computing can significantly decrease the time it takes for data to travel from the device to the cloud and back, resulting in faster response times and improved overall performance \cite{ref8}. This is particularly important for applications that require instantaneous feedback, such as virtual reality, online gaming, and autonomous vehicles \cite{ref9}.

In addition to reducing latency, edge computing also improves real-time processing capabilities \cite{ref2}. By analyzing data in real-time, edge computing can enable faster decision-making, improved predictive maintenance, and enhanced quality control \cite{ref10}. This is particularly important for industries such as manufacturing, healthcare, and finance, where real-time processing can have a significant impact on productivity, safety, and profitability \cite{ref11}. Furthermore, edge computing enhances the overall user experience by providing faster and more reliable access to data and applications \cite{ref1}. By reducing latency and improving real-time processing, edge computing can enable more responsive and interactive applications, which can lead to increased user engagement and satisfaction \cite{ref12}.

The economic and environmental benefits of edge computing are also significant \cite{ref13}. By reducing the need for data to be transmitted to the cloud, edge computing can decrease the amount of energy consumed by data centers and networks, resulting in lower greenhouse gas emissions and operating costs \cite{ref1}. This is particularly important for companies and organizations that are looking to reduce their environmental footprint and improve their sustainability \cite{ref14}. Moreover, edge computing has the potential to enable new business models and revenue streams, particularly in industries such as IoT, smart cities, and industrial automation \cite{ref15}. By providing faster and more reliable access to data and applications, edge computing can enable new services and applications that can generate new revenue streams and improve business productivity \cite{ref1}.

Moreover, edge computing can also improve data security and privacy by reducing the amount of data that needs to be transmitted to the cloud \cite{ref16}. By processing data closer to the source, edge computing can minimize the risk of data breaches and cyber-attacks, which can have significant financial and reputational consequences \cite{ref17}. This is particularly important for industries such as finance, healthcare, and government, where data security and privacy are critical for maintaining public trust and confidence \cite{ref18}. As the field of edge computing continues to evolve, it is likely to have a significant impact on various industries and applications, leading to increased productivity, safety, and profitability. The following section will provide a more in-depth look at the current state of edge computing research, highlighting the key challenges and opportunities in this field \cite{ref2}.

In conclusion, the importance and benefits of edge computing are clear, and its impact is being felt across various industries and applications. By reducing latency, improving real-time processing, and enhancing the overall user experience, edge computing can enable faster and more reliable access to data and applications, which can lead to increased productivity, safety, and profitability \cite{ref2}. As the amount of data being generated continues to grow, edge computing will play an increasingly important role in enabling faster and more reliable access to data and applications, and its benefits will be felt across various industries and applications \cite{ref1}.

\subsection{Key Papers and Research in Edge Computing}

The field of edge computing has undergone significant growth, with a plethora of research papers and studies contributing to its evolution. A seminal paper that defines the landscape of edge computing is \cite{ref2}, which provides a comprehensive overview of the current state of edge computing research, emphasizing its role in reducing latency and improving real-time processing capabilities. This is crucial for applications that require rapid data processing and analysis, such as those discussed in the previous section, including virtual reality, online gaming, and autonomous vehicles \cite{ref9}.

Building on this foundation, other notable papers have explored specific aspects of edge computing. For instance, \cite{ref19} delves into the applications of edge computing in wireless networks, addressing the challenges of enabling ultra-reliable and low-latency edge computing services for mission-critical applications. Furthermore, \cite{ref20} presents a novel multi-tier edge computing framework designed to optimize latency, reduce the probability of failure, and optimize cost, taking into account the sporadic failure of personally owned devices and changing network conditions.

The importance of trustworthiness in edge computing is also a significant area of research, as highlighted in \cite{ref21}, which investigates the use of edge computing for cyber-physical systems with a special emphasis on trustworthiness. Additionally, \cite{ref16} discusses the security threats and vulnerabilities emerging from the edge network architecture, spanning from the hardware layer to the system layer. This security aspect is particularly relevant given the potential of edge computing to improve data security and privacy by reducing the amount of data that needs to be transmitted to the cloud, as discussed earlier \cite{ref16}.

The concept of edge computing is closely intertwined with the idea of pushing artificial intelligence to the wireless network edge, a topic explored in \cite{ref22}. In the context of mobile edge computing, \cite{ref23} provides a novel framework for MEC-enabled heterogeneous networks, composed of multi-tier networks with access points that have different transmission power and computing capabilities. This aligns with the economic and environmental benefits of edge computing, such as reducing the need for data to be transmitted to the cloud, thereby decreasing energy consumption and greenhouse gas emissions \cite{ref1}.

Moreover, edge computing has diverse applications across various fields. For example, \cite{ref24} discusses key considerations for edge deployments in B5G/6G networks, including edge architecture, server location and capacity, user density, and security. The use of edge computing in smart cities is analyzed in \cite{ref25}, which critically reviews the state-of-the-art literature focusing on edge computing applications in smart cities. Additionally, \cite{ref26} presents a survey on the existing research in UAV-enabled edge computing from the resource management perspective, highlighting the versatility of edge computing.

Finally, \cite{ref27} provides a comprehensive survey of open-source edge computing simulators and emulators, emphasizing the convergence of computing and networking paradigms. \cite{ref28} offers a systematic literature review of edge AI, discussing the existing research, recent advancements, and future research directions. These studies collectively underscore the rapid evolution of edge computing and its potential to transform various aspects of technology and society.

In conclusion, the field of edge computing is characterized by its rapid growth and the diverse research efforts aimed at exploring its potential. The papers discussed in this section not only highlight the technical benefits of edge computing, such as reduced latency and improved real-time processing, but also its applications in wireless networks, IoT devices, smart cities, and beyond. As edge computing continues to evolve, addressing challenges related to trustworthiness, security, and privacy will be crucial for its widespread adoption and for realizing its full potential to enable new business models, revenue streams, and innovations, as will be further explored in the subsequent sections.

\section{Edge Computing Architecture}

\subsection{Components of Edge Computing Architecture}

Edge computing architecture is a distributed computing paradigm that brings computation and data storage closer to the source of the data, reducing latency and improving real-time processing capabilities \cite{ref2}. This architecture is foundational to the development and deployment of edge computing applications, which will be discussed in subsequent sections. The basic components of edge computing architecture include edge devices, edge servers, and cloudlets, which interact with each other to provide a seamless computing experience.

Edge devices are the endpoints of the edge computing architecture, such as smartphones, smart home devices, and industrial sensors \cite{ref3}. These devices generate vast amounts of data, which are processed and analyzed at the edge of the network, reducing the need for data transmission to the cloud or central servers. Edge devices can also act as edge servers, providing computing resources to other devices in the network. The data generated by these devices will be processed and analyzed using various edge computing frameworks and platforms, such as those discussed in the following section.

Edge servers, also known as edge nodes or fog nodes, are the intermediate components of the edge computing architecture \cite{ref29}. They are typically located at the edge of the network, such as in cell towers, routers, or switches, and provide computing resources, storage, and networking capabilities to edge devices. Edge servers can be used to process data from multiple edge devices, reducing the latency and bandwidth requirements of the network. The interaction between edge devices and edge servers is critical to the functioning of the edge computing architecture, and will be further enabled by the edge computing frameworks and platforms discussed later.

Cloudlets are small-scale cloud data centers that are located at the edge of the network, typically within a few miles of the edge devices \cite{ref1}. Cloudlets provide a centralized computing resource for edge devices, offering a range of services, including computing, storage, and networking. Cloudlets can be used to process data from multiple edge devices, providing a more scalable and flexible computing resource than edge servers. The use of cloudlets and edge servers will be further discussed in the context of edge computing frameworks and platforms.

The interaction between edge devices, edge servers, and cloudlets is critical to the functioning of the edge computing architecture \cite{ref30}. Edge devices can offload computing tasks to edge servers or cloudlets, which process the data and return the results to the edge devices. Edge servers can also act as intermediaries between edge devices and cloudlets, routing data and computing tasks between the two. This interaction will be further enabled by the development of edge computing frameworks and platforms, which will provide a set of tools, libraries, and APIs for developing, deploying, and managing edge computing applications.

The edge computing architecture also includes a range of other components, such as networking equipment, storage systems, and management software \cite{ref31}. Networking equipment, such as routers and switches, is used to connect edge devices, edge servers, and cloudlets, and to route data between them. Storage systems, such as hard drives and solid-state drives, are used to store data at the edge of the network, reducing the need for data transmission to the cloud or central servers. Management software is used to manage the edge computing architecture, including monitoring and controlling edge devices, edge servers, and cloudlets, and managing data and computing resources. These components will be further discussed in the context of edge computing frameworks and platforms.

In addition to these components, the edge computing architecture also includes a range of enabling technologies, such as containerization, virtualization, and software-defined networking \cite{ref32}. Containerization and virtualization are used to provide a flexible and scalable computing resource, allowing multiple applications and services to run on a single edge server or cloudlet. Software-defined networking is used to manage and control the networking equipment, routing data and computing tasks between edge devices, edge servers, and cloudlets. These enabling technologies will be further discussed in the context of edge computing frameworks and platforms.

The edge computing architecture has a range of benefits, including reduced latency, improved real-time processing capabilities, and increased scalability and flexibility \cite{ref1}. By processing data at the edge of the network, edge computing reduces the need for data transmission to the cloud or central servers, reducing latency and improving real-time processing capabilities. The edge computing architecture also provides a range of services, including computing, storage, and networking, which can be scaled up or down as needed, providing increased scalability and flexibility. These benefits will be further discussed in the context of edge computing frameworks and platforms.

However, the edge computing architecture also has a range of challenges, including security, management, and scalability \cite{ref16}. The edge computing architecture is more complex and distributed than traditional computing architectures, making it more difficult to manage and secure. The edge computing architecture also requires a range of new technologies and standards, including containerization, virtualization, and software-defined networking, which can be complex and difficult to implement. These challenges will be further addressed in the context of edge computing frameworks and platforms.

In conclusion, the edge computing architecture is a distributed computing paradigm that brings computation and data storage closer to the source of the data, reducing latency and improving real-time processing capabilities \cite{ref2}. The basic components of edge computing architecture include edge devices, edge servers, and cloudlets, which interact with each other to provide a seamless computing experience. The edge computing architecture has a range of benefits, including reduced latency, improved real-time processing capabilities, and increased scalability and flexibility, but also has a range of challenges, including security, management, and scalability. The development of edge computing frameworks and platforms will be critical to addressing these challenges and realizing the benefits of edge computing.

\subsection{Edge Computing Frameworks and Platforms}

Edge computing frameworks and platforms are crucial for the development and deployment of edge computing applications, as they provide a set of tools, libraries, and APIs that enable developers to create, deploy, and manage edge computing applications. Building on the edge computing architecture discussed in the previous section, these frameworks and platforms play a vital role in realizing the benefits of edge computing, including reduced latency, improved real-time processing capabilities, and increased scalability and flexibility \cite{ref1}.

One of the popular edge computing frameworks is EdgeFlow \cite{ref30}. EdgeFlow is an open-source, multi-layer edge computing framework that enables the development and deployment of edge computing applications. It provides a set of tools and APIs for data flow processing, computation offloading, and resource allocation. EdgeFlow is designed to support various edge computing scenarios, including IoT, smart cities, and industrial automation. Another notable edge computing framework is FogBus2 \cite{ref33}. FogBus2 is a lightweight, distributed container-based framework that enables the development and deployment of edge computing applications. It provides a set of tools and APIs for resource management, computation offloading, and data processing. FogBus2 is designed to support various edge computing scenarios, including IoT, smart cities, and industrial automation.

In addition to EdgeFlow and FogBus2, there are several other edge computing frameworks and platforms available, including Fog Function \cite{ref34} and Edge-PRUNE \cite{ref35}. These frameworks and platforms provide a set of tools, libraries, and APIs for developing, deploying, and managing edge computing applications. By leveraging these frameworks and platforms, developers can create edge computing applications that can process data in real-time, reduce latency, and improve overall system efficiency. The features and advantages of these frameworks and platforms will be essential in addressing the challenges and open issues in edge computing architecture, which will be discussed in the following section.

Edge computing frameworks and platforms provide several advantages, including reduced latency, improved real-time processing, and enhanced security. By processing data at the edge, edge computing frameworks and platforms can reduce the latency associated with cloud computing. Additionally, edge computing frameworks and platforms can provide real-time processing capabilities, which are essential for applications that require immediate processing and response. Edge computing frameworks and platforms also provide enhanced security features, including data encryption, authentication, and access control. By processing data at the edge, edge computing frameworks and platforms can reduce the risk of data breaches and cyber-attacks.

In terms of features, edge computing frameworks and platforms provide a set of tools and APIs for developing, deploying, and managing edge computing applications. These features include data flow processing, computation offloading, resource allocation, and security management. Edge computing frameworks and platforms are widely used in various industries, including IoT, smart cities, industrial automation, and healthcare. In IoT, edge computing frameworks and platforms are used to develop and deploy applications that require real-time processing and low latency, such as smart home automation and industrial control systems. In smart cities, edge computing frameworks and platforms are used to develop and deploy applications that require real-time processing and low latency, such as intelligent transportation systems and smart energy management.

In industrial automation, edge computing frameworks and platforms are used to develop and deploy applications that require real-time processing and low latency, such as predictive maintenance and quality control. In healthcare, edge computing frameworks and platforms are used to develop and deploy applications that require real-time processing and low latency, such as medical imaging analysis and patient monitoring. As edge computing continues to evolve, we can expect to see more innovative frameworks and platforms emerge, such as \cite{ref34} and \cite{ref35}, which will play a critical role in addressing the challenges and open issues in edge computing architecture.

In conclusion, edge computing frameworks and platforms are essential for the development and deployment of edge computing applications. These frameworks and platforms provide a set of tools, libraries, and APIs for developing, deploying, and managing edge computing applications. Edge computing frameworks and platforms provide several advantages, including reduced latency, improved real-time processing, and enhanced security. They are widely used in various industries, including IoT, smart cities, industrial automation, and healthcare. The discussion of edge computing frameworks and platforms in this section provides a foundation for understanding the challenges and open issues in edge computing architecture, which will be addressed in the following section.

\subsection{Challenges and Open Issues in Edge Computing Architecture}

The edge computing architecture is a complex system that faces numerous challenges and open issues, which can impact the development and deployment of edge computing applications. As discussed in the previous section, edge computing frameworks and platforms provide a set of tools, libraries, and APIs for developing, deploying, and managing edge computing applications. However, the architecture of these systems is crucial to ensuring their scalability, reliability, security, and efficiency. One of the primary concerns in edge computing architecture is scalability, as the number of edge devices and the amount of data they generate continue to grow exponentially \cite{ref2}. To address this issue, researchers have proposed various solutions, such as distributed edge computing architectures and edge-cloud collaborations \cite{ref36}.

These solutions can help to improve the scalability of edge computing systems, but they also introduce new challenges, such as managing the complexity of distributed systems and ensuring seamless communication between edge devices and the cloud. Another significant challenge in edge computing architecture is reliability, as edge devices are often deployed in remote or hard-to-reach locations, making it difficult to maintain and repair them \cite{ref18}. Furthermore, edge devices are prone to failures due to power outages, hardware faults, or software bugs, which can lead to service disruptions and data loss \cite{ref37}. To mitigate these risks, researchers have proposed various reliability mechanisms, such as redundancy, fault tolerance, and self-healing systems \cite{ref38}.

In addition to scalability and reliability, security is another critical concern in edge computing architecture, as edge devices are vulnerable to various types of attacks, including data breaches, denial-of-service attacks, and malware infections \cite{ref16}. Moreover, the decentralized nature of edge computing makes it challenging to ensure the security and integrity of data across the entire system \cite{ref39}. To address these concerns, researchers have proposed various security solutions, such as encryption, access control, and intrusion detection systems \cite{ref40}.

The challenges in edge computing architecture also extend to resource management and optimization, as edge devices often have limited resources, such as computing power, memory, and energy, which must be carefully managed to ensure efficient operation \cite{ref41}. Researchers have proposed various resource management techniques, such as dynamic resource allocation, predictive scheduling, and energy-aware resource allocation \cite{ref42}. Finally, edge computing architecture faces challenges related to standardization and interoperability, as the lack of standardization in edge computing makes it difficult to ensure seamless communication and cooperation between different edge devices and systems \cite{ref1}.

In conclusion, the edge computing architecture faces numerous challenges and open issues related to scalability, reliability, security, resource management, and standardization, which can impact the development and deployment of edge computing applications. Addressing these challenges requires a comprehensive and multidisciplinary approach that involves researchers, developers, and industry stakeholders. By proposing innovative solutions and frameworks, such as distributed edge computing architectures, reliability mechanisms, security solutions, resource management techniques, and standardization efforts, researchers can help mitigate the risks and challenges associated with edge computing and unlock its full potential for various applications and use cases \cite{ref43}. Future research directions may include exploring new edge computing architectures, developing more efficient and effective resource management techniques, and investigating the application of emerging technologies, such as blockchain and artificial intelligence, to edge computing \cite{ref44}.

\section{Applications of Edge Computing}

\subsection{IoT and Smart Cities Applications}

Edge computing has been increasingly applied in various Internet of Things (IoT) and smart city applications, including urban sound classification, smart traffic management, and smart energy management. In urban sound classification, edge computing can be used to analyze and classify different types of sounds in real-time, such as traffic noise, construction noise, and emergency vehicle sirens \cite{ref45}. This can help cities to better understand and manage noise pollution, and to create more livable and sustainable environments for their citizens. By leveraging edge computing, cities can process and analyze large amounts of sensor data in real-time, enabling more efficient and effective management of urban infrastructure.

Smart traffic management is another area where edge computing can be applied to improve the efficiency and safety of transportation systems \cite{ref46}. By analyzing real-time traffic data and sensor inputs, edge computing can help to optimize traffic signal control, reduce congestion, and minimize the risk of accidents. For example, edge computing can be used to analyze data from traffic cameras and sensors to detect incidents and anomalies, and to adjust traffic signal timing accordingly. This can lead to reduced travel times, improved air quality, and enhanced public safety.

In addition to traffic management, edge computing can also be used in smart energy management to optimize energy consumption and reduce waste \cite{ref47}. Edge computing can be used to analyze real-time energy usage data from buildings and homes, and to provide personalized energy usage recommendations to consumers. By leveraging machine learning algorithms and real-time data analytics, edge computing can help to identify areas of inefficiency and optimize energy consumption patterns.

The use of edge computing in IoT and smart city applications has several benefits, including improved real-time processing and analysis, reduced latency, and enhanced security and privacy \cite{ref3}. Edge computing can also help to reduce the amount of data that needs to be transmitted to the cloud, which can help to reduce bandwidth costs and improve network efficiency. Furthermore, edge computing can enable more efficient and effective management of smart city infrastructure, leading to improved quality of life for citizens and enhanced economic competitiveness.

However, the use of edge computing in IoT and smart city applications also poses several challenges, including the need for specialized hardware and software, the complexity of managing and orchestrating edge computing resources, and the need for robust security and privacy measures \cite{ref25}. To address these challenges, researchers and developers are exploring new technologies and techniques, such as fog computing, mobile edge computing, and edge-based machine learning \cite{ref48}. By leveraging these technologies, cities can create more efficient, secure, and sustainable smart city systems that improve the quality of life for citizens.

In conclusion, edge computing has the potential to play a significant role in various IoT and smart city applications, including urban sound classification, smart traffic management, and smart energy management. By analyzing and processing data in real-time, edge computing can help to improve the efficiency, safety, and sustainability of smart cities, and to enhance the overall quality of life for citizens. As the technology continues to evolve and improve, we can expect to see even more innovative applications of edge computing in IoT and smart city contexts, such as smart grids, smart transportation systems, and smart public safety systems \cite{ref49}. The integration of edge computing with other emerging technologies, such as artificial intelligence, blockchain, and the Internet of Things, can help to create more efficient, secure, and sustainable smart cities \cite{ref50}. Overall, edge computing has the potential to play a significant role in the development of smart cities, and to help create more efficient, safe, and sustainable urban environments.

\subsection{Healthcare and Gaming Applications}

Edge computing has numerous applications in various fields, including healthcare and gaming. In healthcare, edge computing can be used for real-time health monitoring, medical imaging analysis, and other applications that require low latency and high bandwidth. For instance, \cite{ref51} discusses the use of edge devices in diagnostics, which can reside within or adjacent to imaging tools such as digital X-ray, CT, MRI, or ultrasound equipment. These devices are either CPUs or GPUs with advanced processing deep and machine learning (artificial intelligence) algorithms that assist in classification and triage solutions to flag studies as either normal or abnormal. This application of edge computing in healthcare is particularly significant, as it enables real-time processing and analysis of medical images, leading to faster and more accurate diagnosis and treatment.

In addition to medical imaging analysis, edge computing can also be used in healthcare for remote patient monitoring. \cite{ref52} proposes a real-time mobile health system for monitoring elderly patients from indoor or outdoor environments. The system uses a bio-signal sensor worn by the patient and a Smartphone as a central node. The sensor data is collected and transmitted to the intelligent server through GPRS/UMTS to be analyzed. The prototype (UMHMSE) monitors the elderly mobility, location, and vital signs such as Sp02 and Heart Rate. Remote users (family and medical personnel) might have a real-time access to the collected information through a web application. This application of edge computing in healthcare highlights the potential of edge computing to improve patient care and outcomes, particularly in remote or resource-constrained areas.

Edge computing also has applications in gaming, particularly in immersive gaming experiences. \cite{ref53} proposes a closed-loop architecture integrating the Network Data Analytics Function (NWDAF) and UPF to estimate user latency and enhance the 5G control plane by making it latency-aware. The results show that embedding an artificial intelligence model within NWDAF enables game classification and opens new avenues for mobile edge gaming research. Furthermore, edge computing can be used in gaming for real-time rendering and processing, as well as predictive analytics and personalized gaming experiences. \cite{ref54} presents an XR system for enabling flexible edge-assisted XR, which uses a cascaded WLAN and a cellular network to collect physiological and environmental data in-home and transfer them to a remote hospital.

The use of edge computing in healthcare and gaming has the potential to transform these industries, enabling new use cases and applications that were not possible with traditional cloud computing. As edge computing continues to evolve, we can expect to see even more innovative applications in these fields. The integration of edge computing with other technologies such as artificial intelligence, IoT, and 5G networks can further enhance its applications in healthcare and gaming. For instance, \cite{ref55} discusses the use of AI at the edge, which can enable real-time processing and analysis of data, and can improve the accuracy and responsiveness of various applications. Similarly, \cite{ref56} proposes an IoT-enabled low-cost fog computing system with online machine learning for accurate and low-latency heart monitoring in rural healthcare settings.

However, the use of edge computing in healthcare and gaming also raises several challenges and concerns, such as security, privacy, and scalability. \cite{ref57} discusses the security and privacy issues in modern healthcare systems, and proposes several solutions to address these challenges. Similarly, \cite{ref58} proposes a privacy-preserving healthcare IoT security management system, which can protect the privacy and security of patients' data. To address these challenges, it is essential to develop context-aware approaches and frameworks that can ensure the efficient, accurate, and responsive use of edge computing in healthcare and gaming. \cite{ref59} discusses the opportunities and challenges of edge computing in smart health, and proposes several context-aware approaches to address these challenges. \cite{ref60} proposes an edge learning as a service framework for knowledge-centric connected healthcare, which can enable the efficient, accurate, and responsive use of edge computing in healthcare.

In conclusion, edge computing has numerous applications in healthcare and gaming, including real-time health monitoring, medical imaging analysis, immersive gaming experiences, and predictive analytics. The use of edge computing in these fields can improve the efficiency, accuracy, and responsiveness of various applications, and can enable new use cases that were not possible with traditional cloud computing. As edge computing continues to evolve, it is likely to have a significant impact on the development of smart healthcare and gaming systems, and its integration with other technologies such as AI, IoT, and 5G networks can further enhance its applications in these fields. Ultimately, the future of edge computing in healthcare and gaming is exciting and challenging, with significant opportunities for growth and innovation, and its potential to transform these industries is vast and promising.

\subsection{Future Directions and Challenges of Edge Computing Applications}

As edge computing continues to evolve and gain traction, several future directions and challenges emerge that will shape the development and adoption of edge computing applications. Building on the numerous applications of edge computing in healthcare and gaming, as well as its integration with other technologies such as artificial intelligence, IoT, and 5G networks, it is essential to address the challenges that will enable the efficient, accurate, and responsive use of edge computing. One of the primary challenges is scalability, as edge computing systems need to be able to handle increasing amounts of data and traffic while maintaining low latency and high performance \cite{ref2}. To address this challenge, researchers and developers are exploring new architectures and technologies, such as fog computing and mobile edge computing, that can provide more scalable and flexible edge computing solutions \cite{ref19}.

Another significant challenge is reliability, as edge computing systems are often deployed in remote or hard-to-reach locations, making it difficult to maintain and repair them \cite{ref18}. To improve reliability, edge computing systems need to be designed with redundancy and failover capabilities, as well as advanced monitoring and management tools to quickly detect and respond to faults and errors \cite{ref61}. Additionally, edge computing systems need to be secured against cyber threats, which is a significant challenge due to the distributed and heterogeneous nature of edge computing environments \cite{ref16}. Standardization is also a critical challenge for edge computing applications, as different vendors and organizations are developing their own proprietary edge computing solutions, which can lead to interoperability issues and fragmentation \cite{ref1}. To address this challenge, industry-wide standards and frameworks are needed to ensure interoperability and compatibility between different edge computing solutions \cite{ref62}.

In terms of future directions, edge computing is expected to play a critical role in the development of emerging technologies, such as artificial intelligence, Internet of Things (IoT), and 5G networks \cite{ref63}. Edge computing can provide the low-latency and high-performance processing needed to support these technologies, and can also help to reduce the amount of data that needs to be transmitted to the cloud, improving overall efficiency and reducing costs \cite{ref24}. The development of edge-based machine learning and artificial intelligence is also a promising area, which can provide real-time insights and decision-making capabilities at the edge of the network \cite{ref43}. This can be particularly useful in applications, such as smart cities, industrial automation, and healthcare, where real-time processing and decision-making are critical \cite{ref64}. However, the development of edge-based machine learning and artificial intelligence also raises significant challenges, such as the need for large amounts of training data, the complexity of machine learning algorithms, and the need for specialized hardware and software \cite{ref6}. To address these challenges, researchers and developers are exploring new techniques, such as federated learning and transfer learning, that can provide more efficient and effective edge-based machine learning and artificial intelligence solutions \cite{ref36}.

Ultimately, the future of edge computing applications is exciting and challenging, with significant opportunities for growth and innovation. By addressing the challenges of scalability, reliability, standardization, and the development of edge-based machine learning and artificial intelligence, edge computing can provide the low-latency, high-performance, and real-time processing needed to support emerging technologies and applications, and can help to create a more efficient, effective, and sustainable computing paradigm \cite{ref65}. As edge computing continues to evolve, it is likely to have a profound impact on various industries and applications, and its potential to transform the way we live and work is vast and promising.

\section{Edge Computing Techniques and Technologies}

\subsection{Computation Offloading and Resource Allocation}

Computation offloading and resource allocation are crucial components of edge computing, as they enable the efficient distribution of computational tasks between edge devices and the cloud. The primary goal of computation offloading is to minimize latency and reduce the energy consumption of edge devices, while resource allocation aims to optimize the utilization of available resources, such as computing power, memory, and bandwidth. In edge computing, the integration of computation offloading and resource allocation is essential to support various applications, including those that require real-time processing, low latency, and high throughput.

One of the key challenges in edge computing is the limited resources available on edge devices, such as smartphones, smart home devices, and vehicles. These devices often have limited computing power, memory, and storage, making it difficult to execute computationally intensive tasks, such as video processing, natural language processing, and machine learning. Computation offloading addresses this challenge by offloading tasks to more powerful edge servers or the cloud, reducing the computational burden on edge devices and improving overall system performance \cite{ref66}. This approach enables edge devices to conserve energy, reduce latency, and improve overall system efficiency.

To optimize computation offloading and resource allocation in edge computing, various approaches have been proposed. One approach is to use game theory to model the interaction between edge devices and edge servers, and to develop algorithms that optimize the offloading decision and resource allocation \cite{ref67}. Another approach is to use machine learning and deep learning techniques to predict the computational requirements of tasks and optimize the offloading decision and resource allocation accordingly \cite{ref68}. These approaches enable edge computing systems to adapt to changing environmental conditions, optimize resource utilization, and improve overall system performance.

Resource allocation is also a critical component of edge computing, as it determines how available resources are allocated to different tasks and applications. Various resource allocation algorithms have been proposed, including those based on convex optimization, dynamic programming, and reinforcement learning \cite{ref69}. These algorithms aim to optimize the allocation of resources, such as computing power, memory, and bandwidth, to minimize latency and improve overall system performance. By optimizing resource allocation, edge computing systems can support multiple applications and services, while ensuring low latency, high throughput, and efficient resource utilization.

In addition to computation offloading and resource allocation, other techniques have been proposed to optimize edge computing performance, including data compression, caching, and content delivery networks \cite{ref70}. These techniques aim to reduce the amount of data that needs to be transmitted between edge devices and edge servers, reducing latency and improving overall system performance. By combining computation offloading, resource allocation, and these techniques, edge computing systems can achieve optimal performance, low latency, and high efficiency.

The emergence of new edge computing architectures, such as fog computing and cloudlet-based edge computing, has also led to the development of new computation offloading and resource allocation algorithms \cite{ref71}. These algorithms aim to optimize the offloading decision and resource allocation in these new architectures, taking into account the unique characteristics of each architecture. By developing new algorithms and techniques, researchers can address the challenges associated with computation offloading and resource allocation in edge computing, and enable the widespread adoption of edge computing technologies.

In summary, computation offloading and resource allocation are critical components of edge computing, and their optimization is essential to improve the performance of edge computing systems. Various approaches and algorithms have been proposed to optimize computation offloading and resource allocation, including those based on game theory, machine learning, and optimization theory \cite{ref72}. The development of new edge computing architectures has also led to the development of new computation offloading and resource allocation algorithms, highlighting the need for continued research and innovation in this area \cite{ref73}. By optimizing computation offloading and resource allocation, edge computing systems can achieve low latency, high efficiency, and optimal performance, enabling the widespread adoption of edge computing technologies in various applications and services.

The optimization of computation offloading and resource allocation in edge computing is a complex problem that requires careful consideration of various factors, including the computational requirements of tasks, the availability of resources, and the latency and energy consumption constraints of edge devices. To address this challenge, researchers have proposed various algorithms and techniques, including those based on machine learning, game theory, and optimization theory \cite{ref74}. By developing new algorithms and techniques, researchers can address the challenges associated with computation offloading and resource allocation in edge computing, and enable the widespread adoption of edge computing technologies in various applications and services. As edge computing continues to evolve, the optimization of computation offloading and resource allocation will play a critical role in enabling the development of efficient, low-latency, and high-performance edge computing systems.

\subsection{Edge-Based Machine Learning and Artificial Intelligence}

Edge-based machine learning and artificial intelligence are crucial components of edge computing, enabling the development of intelligent edge systems that can process and analyze data in real-time. The integration of machine learning and artificial intelligence with edge computing has given rise to various applications, including edge-based deep learning, reinforcement learning, and federated learning \cite{ref75}. These techniques have the potential to revolutionize the way data is processed and analyzed at the edge, enabling faster and more accurate decision-making.

As discussed earlier, computation offloading and resource allocation are essential components of edge computing, and they play a critical role in enabling the efficient deployment of machine learning and artificial intelligence models on edge devices. The use of edge-based machine learning and artificial intelligence can significantly improve the performance of edge computing systems, particularly in applications that require real-time processing and analysis of data. For instance, edge-based deep learning can be used to perform tasks such as image recognition, natural language processing, and predictive analytics \cite{ref7}. The use of deep learning at the edge enables real-time processing and analysis of data, reducing the latency and bandwidth requirements associated with transmitting data to the cloud for processing.

Reinforcement learning is another important application of machine learning in edge computing, where agents learn to make decisions based on rewards or penalties received from the environment \cite{ref76}. Reinforcement learning has been used in various edge computing applications, including robotics, autonomous vehicles, and smart homes. The use of reinforcement learning at the edge enables real-time decision-making and adaptation to changing environmental conditions, reducing the need for human intervention and improving overall system efficiency. Furthermore, federated learning is a distributed machine learning approach that enables multiple edge devices to collaboratively train a shared model while keeping their data private \cite{ref77}.

The application of edge-based machine learning and artificial intelligence has numerous benefits, including improved real-time processing and analysis, reduced latency and bandwidth requirements, and enhanced privacy and security \cite{ref78}. The use of edge-based machine learning and artificial intelligence also enables the development of autonomous systems that can operate independently, reducing the need for human intervention and improving overall system efficiency. However, the deployment of edge-based machine learning and artificial intelligence also poses significant challenges, including the need for specialized hardware and software, as well as the requirement for large amounts of training data \cite{ref7}.

To address these challenges, researchers have proposed various techniques, including model pruning, knowledge distillation, and federated learning \cite{ref77}. Model pruning involves reducing the complexity of deep neural networks to enable their deployment on resource-constrained edge devices. Knowledge distillation involves transferring knowledge from a large pre-trained model to a smaller model, enabling the deployment of accurate models on edge devices. These techniques can be used in conjunction with optimization techniques, such as convex optimization, reinforcement learning, and game theory, to further improve the performance of edge computing systems.

In conclusion, edge-based machine learning and artificial intelligence are crucial components of edge computing, enabling the development of intelligent edge systems that can process and analyze data in real-time. The integration of machine learning and artificial intelligence with edge computing has given rise to various applications, including edge-based deep learning, reinforcement learning, and federated learning \cite{ref75}. While there are significant benefits to the use of edge-based machine learning and artificial intelligence, there are also challenges that need to be addressed, including the need for specialized hardware and software, as well as the requirement for large amounts of training data \cite{ref6}. The use of edge-based machine learning and artificial intelligence will require significant advances in areas such as computer vision, natural language processing, and predictive analytics \cite{ref79}, and will have significant implications for the future of edge computing, particularly in areas such as robotics, autonomous vehicles, and smart homes \cite{ref80}.

\subsection{Optimization Techniques for Edge Computing}

Optimization techniques are essential in edge computing as they enable the efficient allocation of resources, reduction of latency, and improvement of overall system performance. Building on the concepts of edge-based machine learning and artificial intelligence discussed earlier, various optimization techniques have been applied to edge computing, including convex optimization, reinforcement learning, and game theory. These techniques are crucial in optimizing the performance of edge computing systems, particularly in applications that involve real-time processing and analysis of data.

Convex optimization is a powerful tool for optimizing edge computing systems, and it has been used to optimize resource allocation, reduce latency, and improve system performance. For example, \cite{ref81} presents a comprehensive overview of conic optimization theory, including convexification techniques and numerical algorithms. Convex optimization has also been used to optimize the placement of edge servers, as shown in \cite{ref82}. This paper formulates the resource dimensioning problem as a convex optimization problem and provides numerical results to demonstrate the effectiveness of the proposed approach.

Reinforcement learning is another optimization technique that has been widely applied to edge computing, enabling edge devices to learn from their environment and make decisions to optimize system performance. For example, \cite{ref83} proposes a novel reinforcement learning algorithm for edge computing, which achieves significantly improved performance with reduced computational resources. Reinforcement learning has also been used to optimize task offloading in edge computing, as shown in \cite{ref84}. This paper formulates the task offloading problem as a reinforcement learning problem and proposes a deep reinforcement learning approach to solve it.

Game theory is also a useful optimization technique for edge computing, enabling edge devices to make decisions that optimize system performance while taking into account the actions of other devices. For example, \cite{ref85} proposes a game-theoretic approach to optimize collaborative edge computing. This paper formulates the collaborative edge computing problem as a game-theoretic problem and proposes a deep reinforcement learning approach to solve it. Game theory has also been used to optimize resource allocation in edge computing, as shown in \cite{ref66}. This paper formulates the resource allocation problem as a game-theoretic problem and proposes a convex optimization approach to solve it.

In addition to these optimization techniques, other approaches have also been proposed to optimize edge computing systems, such as federated learning and deep reinforcement learning. For example, \cite{ref86} proposes a deep reinforcement learning approach to optimize function placement in edge computing. This paper formulates the function placement problem as a reinforcement learning problem and proposes a deep reinforcement learning approach to solve it. \cite{ref80} proposes a federated learning approach to optimize edge computing, caching, and communication. This paper formulates the optimization problem as a federated learning problem and proposes a federated learning approach to solve it.

The application of these optimization techniques has significant implications for the future of edge computing, particularly in areas such as robotics, autonomous vehicles, and smart homes. As edge computing continues to evolve, optimization techniques will play an increasingly important role in optimizing system performance and reducing latency. \cite{ref87} presents a comprehensive overview of fast convex optimization techniques for two-layer ReLU networks, which can be applied to edge computing. \cite{ref88} provides a comprehensive tutorial on advanced optimization methods, including convex optimization, linearization techniques, and mixed-integer linear programming, which can be used to optimize edge computing systems.

In conclusion, optimization techniques are crucial in edge computing, and various techniques have been proposed to optimize edge computing systems. Convex optimization, reinforcement learning, and game theory are some of the most widely used optimization techniques in edge computing, and they have been used to optimize resource allocation, reduce latency, and improve system performance. As edge computing continues to evolve, the application of these optimization techniques will be essential in optimizing system performance and reducing latency, enabling the development of more efficient and effective edge computing systems.

\section{Security and Privacy in Edge Computing}

\subsection{Security Threats and Vulnerabilities}

5.1 Security Threats and Vulnerabilities

Edge computing, despite its numerous benefits, is not immune to security threats and vulnerabilities. The distributed nature of edge computing, where data is processed at the edge of the network, closer to the source, introduces a unique set of challenges. One of the primary concerns is the risk of data breaches \cite{ref16}, which can occur due to unauthorized access or malicious attacks on edge devices. These breaches can compromise sensitive information, leading to significant financial and reputational losses. As edge computing continues to grow and expand into various industries, the potential attack surface also increases, making it essential to address these security concerns.

The security threats in edge computing can be categorized into several types, including unauthorized access, malicious attacks, and physical attacks. Unauthorized access is a significant threat, as edge devices are more susceptible to physical tampering, which can allow attackers to gain unauthorized access to the system \cite{ref89}. This can be particularly problematic in applications such as intelligent transportation systems, where the consequences of a security breach can be severe. Malicious attacks, including distributed denial-of-service (DDoS) attacks, malware injections, and side-channel attacks, are also major concerns in edge computing \cite{ref90}. These attacks can compromise the integrity of edge devices, leading to a loss of trust in the system.

Furthermore, the use of machine learning models in edge computing introduces additional vulnerabilities, such as model inversion attacks and membership inference attacks \cite{ref17}. The lack of standardization in edge computing also poses significant security risks \cite{ref18}. The heterogeneity of edge devices and the lack of interoperability between different systems can make it challenging to implement robust security measures. Moreover, the limited resources available on edge devices can limit the effectiveness of security solutions, making them more vulnerable to attacks. Physical attacks, such as electromagnetic fault injection (EMFI) attacks, can also compromise the integrity of edge devices \cite{ref91}.

The security threats and vulnerabilities in edge computing can have severe consequences, including financial losses, reputational damage, and compromised safety \cite{ref92}. Therefore, it is essential to develop robust security measures to mitigate these risks. This can include implementing secure authentication and authorization mechanisms, encrypting data both in transit and at rest, and using secure communication protocols \cite{ref93}. Additionally, the use of blockchain technology can provide an additional layer of security in edge computing \cite{ref39}. Blockchain-based solutions can provide a secure and decentralized way to manage data and ensure the integrity of edge devices. As the field of edge computing continues to evolve, it is crucial to address these security concerns and develop innovative solutions to ensure the integrity and trustworthiness of edge computing systems, which will be discussed in the following section on security and privacy solutions.

\subsection{Security and Privacy Solutions}

To address the security and privacy concerns in edge computing, various solutions have been proposed, including encryption, authentication, and access control mechanisms. As discussed in the previous section, edge computing is vulnerable to numerous security threats and vulnerabilities, and it is essential to develop robust security measures to mitigate these risks. Encryption is a crucial technique for protecting data in edge computing systems, and \cite{ref94} presents a scheme that guarantees information-theoretic user data privacy against an eavesdropper with access to a number of edge servers or their corresponding communication links.

In addition to encryption, authentication is another essential security mechanism in edge computing. \cite{ref95} proposes an efficient and lightweight mutual authentication protocol for MEC environments, based on Elliptic Curve Cryptography (ECC), one-way hash functions, and concatenation operations. Access control is also a critical security mechanism in edge computing, and \cite{ref96} proposes an adaptive and efficient authentication and key agreement scheme for Cloud-Edge-Device IoT environments.

Several other solutions have been proposed to address specific security and privacy concerns in edge computing. For example, \cite{ref97} presents a blockchain-powered edge intelligence framework for u-healthcare applications, which provides a secure and privacy-preserving environment for data processing and analysis. \cite{ref98} proposes a secure and fine-grained access control scheme for in-vehicle communications based on edge computing. \cite{ref99} presents a secure coded cooperative computation mechanism for edge computing, which provides both security and computation efficiency guarantees.

Furthermore, \cite{ref100} proposes a reconfigurable security framework for IoT based on edge computing, which utilizes a near-user edge device to simplify key management and offload the computational costs of security algorithms at IoT devices. \cite{ref90} addresses the issue of membership inference leakage in distributed edge intelligence systems, where edge nodes collaboratively generate a global model. \cite{ref101} presents a solution for edge-ready and end-to-end secure protocol extensions, which can efficiently maintain end-to-edge-to-end (E\^{}3) security semantics.

Other notable solutions include \cite{ref102}, which proposes a probabilistic approach for privacy-preserving edge caching, and \cite{ref103}, which designs a smart contract-based decentralized technique to establish the trading trust among users and provide flexible smart contract interfaces to satisfy users. \cite{ref104} proposes a unified framework for private edge computing, which studies two highly-coupled problems: storage allocation and computation design. \cite{ref105} presents an efficient software implementation framework for pairing-based function-hiding inner product encryption (FHIPE) using the recently proposed and widely adopted BLS12-381 pairing-friendly elliptic curve. \cite{ref106} proposes a lightweight blockchain-based access control scheme for integrated edge computing in IoT, which integrates attribute-based encryption (ABE) and attribute-based access control (ABAC) models with attributes as the entry point.

These security and privacy solutions for edge computing have shown effectiveness in protecting edge computing systems from various threats and vulnerabilities. However, the edge computing environment is constantly evolving, and new security and privacy challenges are emerging. Therefore, it is essential to continue researching and developing new security and privacy solutions to address these challenges and ensure the secure and private operation of edge computing systems. The future research directions and open challenges in security and privacy for edge computing will be discussed in the following section, which will explore the technical, regulatory, and human factors challenges that need to be addressed to ensure the widespread adoption of edge computing technologies.

\subsection{Future Research Directions and Open Challenges}

Future research directions and open challenges in security and privacy for edge computing are numerous and multifaceted, building upon the existing security and privacy solutions discussed earlier. One of the primary concerns is the need for more robust security protocols that can effectively protect edge devices and the data they process \cite{ref107}. The current security measures are often inadequate, and the lack of standardization in edge computing security protocols exacerbates the problem \cite{ref92}. Therefore, developing and implementing robust security protocols that can ensure the confidentiality, integrity, and availability of data in edge computing environments is crucial, building on the foundations established by existing solutions such as encryption, authentication, and access control mechanisms.

Another significant challenge is the development of better authentication mechanisms \cite{ref96}. Edge devices often have limited computational resources and power, making it difficult to implement complex authentication protocols. Moreover, the dynamic nature of edge computing environments, with devices constantly joining and leaving the network, requires authentication mechanisms that are not only secure but also efficient and scalable \cite{ref108}. Researchers and developers must work together to design and implement authentication protocols that balance security with the constraints of edge devices, taking into account the lessons learned from existing authentication protocols.

The need for more effective privacy-preserving techniques is also a pressing concern \cite{ref109}. Edge computing involves processing sensitive data from various sources, including IoT devices, smartphones, and other personal devices. Ensuring that this data is protected from unauthorized access and misuse is essential \cite{ref94}. Techniques such as differential privacy, homomorphic encryption, and secure multi-party computation can be explored to develop privacy-preserving edge computing applications \cite{ref110}. However, these techniques often come with significant computational overhead, which can be challenging for resource-constrained edge devices, highlighting the need for innovative solutions that balance privacy and performance.

In addition to these technical challenges, there are also regulatory and standardization challenges that need to be addressed \cite{ref111}. The lack of standardization in edge computing security and privacy can lead to a fragmented market, making it difficult for developers to create secure and interoperable applications. Furthermore, regulatory frameworks for edge computing security and privacy are still in their infancy, and there is a need for clearer guidelines and regulations to ensure that edge computing applications are developed and deployed in a secure and privacy-preserving manner \cite{ref112}. Addressing these challenges will be crucial for the widespread adoption of edge computing technologies.

The emergence of new technologies such as blockchain, artificial intelligence, and the Internet of Things (IoT) also presents both opportunities and challenges for edge computing security and privacy \cite{ref113}. For example, blockchain technology can be used to develop secure and decentralized edge computing applications \cite{ref38}, while artificial intelligence can be used to develop more effective security and privacy mechanisms \cite{ref40}. However, these technologies also introduce new security and privacy risks that need to be addressed, such as the potential for blockchain-based applications to be vulnerable to 51\% attacks or the potential for AI-based systems to be biased or discriminatory, requiring careful consideration and mitigation strategies.

Finally, there is a need for more research on the human factors aspect of edge computing security and privacy \cite{ref114}. Edge computing applications often involve human users, and understanding how users perceive and interact with these applications is crucial for developing effective security and privacy mechanisms. For example, research has shown that users are often unaware of the security and privacy risks associated with edge computing applications, and that they may not take adequate precautions to protect their data \cite{ref115}. Therefore, developing user-centered security and privacy mechanisms that take into account the needs and limitations of human users is essential for ensuring the security and privacy of edge computing applications, and will be an important area of focus in future research.

In conclusion, the future research directions and open challenges in security and privacy for edge computing are numerous and complex, requiring a comprehensive and multidisciplinary approach to address the technical, regulatory, and human factors challenges. By building on existing solutions, addressing emerging challenges, and exploring new technologies and approaches, researchers and developers can work together to ensure that edge computing applications are secure, private, and trustworthy \cite{ref116}, ultimately unlocking the full potential of edge computing technologies.

\section{Resource Management and Optimization}

\subsection{Resource Allocation Techniques}

Resource allocation is a critical component of edge computing, as it directly impacts the performance, efficiency, and scalability of edge-based applications. Various resource allocation techniques have been proposed to optimize the utilization of edge resources, including dynamic resource allocation, predictive scheduling, and energy-aware resource allocation \cite{ref117}. These techniques are essential to ensure that edge computing systems can efficiently manage their resources and provide high-quality services to users.

Dynamic resource allocation is a technique that involves allocating resources dynamically based on the changing demands of edge applications \cite{ref118}. This approach enables edge systems to adapt to varying workloads and optimize resource utilization, which is crucial for applications that require low latency and high throughput. By dynamically allocating resources, edge systems can ensure that resources are not underutilized or overutilized, which can lead to improved efficiency and reduced energy consumption.

Predictive scheduling is another technique that involves predicting the future demands of edge applications and allocating resources accordingly \cite{ref119}. This approach enables edge systems to prepare for upcoming workloads and optimize resource utilization, which can lead to improved performance and reduced energy consumption. Predictive scheduling can be particularly useful in applications where the workload is predictable, such as in smart cities or industrial automation.

Energy-aware resource allocation is a technique that involves allocating resources based on the energy consumption of edge devices \cite{ref120}. This approach enables edge systems to minimize energy consumption and reduce the carbon footprint of edge applications, which is essential for applications that require low power consumption, such as in IoT devices. By allocating resources based on energy consumption, edge systems can ensure that energy-intensive tasks are allocated to devices that have sufficient power resources, which can lead to improved energy efficiency.

These resource allocation techniques have various applications in different edge computing scenarios, including IoT, smart cities, and healthcare \cite{ref67}. For instance, in IoT applications, dynamic resource allocation can be used to allocate resources based on the changing demands of IoT devices \cite{ref121}. In smart cities, predictive scheduling can be used to allocate resources based on the predicted demands of smart city applications \cite{ref122}.

In addition to these techniques, other approaches have been proposed to optimize resource allocation in edge computing, such as joint task offloading scheduling and transmit power allocation \cite{ref123}, and energy-aware resource allocation \cite{ref124}. These approaches can be used to optimize resource allocation in edge computing systems and improve the quality of service (QoS) of edge applications.

Several studies have proposed various frameworks and techniques for resource allocation in edge computing, including \cite{ref125}, \cite{ref126}, \cite{ref127}, \cite{ref128}, \cite{ref129}, \cite{ref130}, \cite{ref131}, \cite{ref132}, \cite{ref133}, \cite{ref134}, \cite{ref135}, \cite{ref136}, \cite{ref137}, \cite{ref138}, \cite{ref139}, \cite{ref140}, \cite{ref141}, \cite{ref142}, \cite{ref143}, \cite{ref144}, \cite{ref145}, \cite{ref146}, \cite{ref147}, \cite{ref148}, \cite{ref149}, \cite{ref150}, \cite{ref151}, \cite{ref152}. These studies demonstrate the importance of resource allocation in edge computing and the need for efficient resource allocation techniques to optimize the performance and energy efficiency of edge computing systems.

In conclusion, resource allocation is a critical component of edge computing, and various techniques have been proposed to optimize the utilization of edge resources. These techniques have various applications in different edge computing scenarios and can improve the QoS of edge applications, minimize energy consumption, and reduce the carbon footprint of edge computing. By optimizing resource allocation, edge computing systems can provide high-quality services to users while minimizing energy consumption and reducing the carbon footprint of edge applications, which is essential for the widespread adoption of edge computing in various industries. The optimization of computation and communication resources, which is discussed in the next section, is closely related to resource allocation and is essential to achieve optimized performance in edge computing systems.

\subsection{Optimization of Computation and Communication Resources}

The optimization of computation and communication resources is a crucial aspect of edge computing, as it directly impacts the performance, energy efficiency, and latency of edge-based applications. Building on the resource allocation techniques discussed in the previous section, the joint optimization of computation and communication resources is essential to minimize energy consumption and latency, and to ensure that edge computing systems can efficiently manage their resources and provide high-quality services to users. In edge computing, computation and communication resources are intertwined, and optimizing one without considering the other can lead to suboptimal performance.

To address this challenge, various approaches can be employed to optimize computation and communication resources. One approach is to use machine learning techniques, such as the federated edge learning framework proposed in \cite{ref153}, which jointly optimizes communication and computation resources to minimize energy consumption. Another approach is to use optimization techniques, such as linear programming or convex optimization, to allocate computation and communication resources, as demonstrated in \cite{ref67}, which formulates a mixed-integer linear programming problem to jointly optimize task offloading scheduling and transmit power allocation in mobile-edge computing systems.

In addition to machine learning and optimization techniques, other approaches, such as game theory and auction theory, can be used to optimize computation and communication resources in edge computing. For instance, \cite{ref69} proposes a game-theoretic approach to allocate computation and communication resources in mobile-edge computing systems. These approaches can be applied to various edge computing scenarios, including real-time applications such as video analytics and autonomous vehicles, which require optimized computation and communication resources to meet stringent latency and throughput requirements.

The optimization of computation and communication resources is closely related to the concept of computation offloading, which involves offloading computation-intensive tasks from edge devices to more powerful edge servers or cloud servers. \cite{ref154} proposes a computation offloading framework that jointly optimizes computation and communication resources to maximize computation rates in wireless powered edge computing systems. Furthermore, the optimization of computation and communication resources is also related to the concept of resource allocation, which involves allocating limited resources, such as bandwidth, computation power, and energy, to different edge devices and applications. \cite{ref155} proposes a resource allocation framework that jointly optimizes computation and communication resources to minimize latency and energy consumption in semantic-aware mobile-edge computing systems.

In conclusion, the optimization of computation and communication resources is a critical aspect of edge computing, and various approaches, such as machine learning, optimization techniques, game theory, and auction theory, can be used to achieve optimized performance. The use of edge computing in real-time applications requires optimized computation and communication resources to meet stringent latency and throughput requirements. By jointly optimizing computation and communication resources, edge computing systems can achieve low latency, high throughput, and energy efficiency, making them suitable for a wide range of applications. As discussed in the following section, the application of edge computing in various industries, such as smart cities, healthcare, and industrial automation, relies on efficient resource management and optimization techniques, including the optimization of computation and communication resources, to ensure low latency, high throughput, and reliable operation, as demonstrated in \cite{ref156} and \cite{ref157}.

\subsection{Case Studies and Applications}

Case studies and applications of resource management and optimization in edge computing are crucial in understanding the practical implications of these techniques, as they demonstrate how the optimization of computation and communication resources, discussed in the previous section, can be applied in real-world scenarios. Various industries, such as smart cities, healthcare, and industrial automation, have benefited from the deployment of edge computing solutions, which rely on efficient resource management and optimization to ensure low latency, high throughput, and reliable operation. For instance, in smart cities, edge computing can be used to manage and analyze data from various sensors and IoT devices, enabling real-time monitoring and decision-making \cite{ref3}. This can lead to improved public safety, traffic management, and energy efficiency, which are all critical aspects of smart city infrastructure.

In the healthcare sector, edge computing can be applied to analyze medical images and patient data in real-time, enabling faster diagnosis and treatment \cite{ref158}. For example, edge-based AI models can be used to detect abnormalities in medical images, such as tumors or fractures, and provide instant feedback to medical professionals. This can significantly improve patient outcomes and reduce the workload of healthcare professionals, highlighting the importance of optimized computation and communication resources in edge computing. The use of edge computing in healthcare also requires careful consideration of resource allocation and optimization, as medical applications often have strict latency and throughput requirements.

Industrial automation is another area where edge computing can have a significant impact, as it enables the analysis of data from machines and sensors in real-time, allowing for predictive maintenance and improved overall equipment effectiveness \cite{ref159}. This can lead to reduced downtime, improved product quality, and increased productivity, all of which are critical aspects of industrial automation. The application of edge computing in industrial automation also relies on efficient resource management and optimization, as the analysis of large amounts of data from machines and sensors requires significant computational resources.

Edge computing can also be applied in other areas, such as transportation systems, where it can be used to analyze data from sensors and cameras to optimize traffic flow and reduce congestion \cite{ref160}. Additionally, edge computing can be used in smart homes and buildings to manage and control various systems, such as lighting, heating, and security, leading to improved energy efficiency and comfort \cite{ref161}. These applications all rely on efficient resource management and optimization, highlighting the importance of developing techniques that can optimize computation and communication resources in edge computing environments.

The use of edge computing in these applications requires efficient resource management and optimization techniques to ensure low latency, high throughput, and reliable operation. Various techniques, such as containerization, orchestration, and edge-based machine learning, can be used to optimize resource allocation and utilization in edge computing environments \cite{ref162}. Real-world examples of edge computing deployments can be found in various industries, demonstrating the potential of edge computing to transform industries and improve efficiency, productivity, and decision-making. For instance, the city of Barcelona has deployed an edge computing solution to manage and analyze data from various sensors and IoT devices, enabling real-time monitoring and decision-making \cite{ref3}. Similarly, the healthcare company, Philips, has developed an edge-based AI solution for medical image analysis, which can detect abnormalities in medical images and provide instant feedback to medical professionals \cite{ref158}.

In industrial automation, companies like Siemens and GE have developed edge computing solutions to analyze data from machines and sensors in real-time, enabling predictive maintenance and improving overall equipment effectiveness \cite{ref159}. These examples demonstrate the potential of edge computing to transform various industries and improve efficiency, productivity, and decision-making, and highlight the importance of continued research into resource management and optimization techniques for edge computing. As the demand for edge computing continues to grow, it is essential to develop efficient resource management and optimization techniques to ensure low latency, high throughput, and reliable operation \cite{ref117}, which will be discussed in more detail in the following section.

\section{Edge Computing and Artificial Intelligence}

\subsection{Introduction to Edge AI}

The concept of Edge AI has emerged as a revolutionary paradigm in the field of artificial intelligence, enabling real-time processing and reducing latency in various applications. The history of Edge AI dates back to the early 2000s, when researchers began exploring the idea of processing data closer to the source, rather than relying on cloud-based computing \cite{ref2}. This shift towards edge computing was driven by the need for low-latency and real-time processing capabilities, which is essential for applications such as autonomous vehicles, smart cities, and industrial automation \cite{ref19}. The importance of Edge AI lies in its ability to provide reduced latency and enable real-time decision-making, making it a crucial component in these applications.

One of the key drivers of Edge AI is the need for reduced latency in real-time applications. Traditional cloud-based computing architectures are often plagued by high latency, which can be detrimental to applications that require immediate processing and decision-making \cite{ref163}. To address this issue, Edge AI processes data closer to the source, thereby reducing the latency and enabling real-time decision-making. For instance, in autonomous vehicles, Edge AI can process sensor data in real-time, enabling the vehicle to make immediate decisions and react to changing environments \cite{ref164}. This capability is critical for ensuring the safety and efficiency of autonomous vehicles.

The emergence of large language models (LLMs) has also played a significant role in the development of Edge AI \cite{ref165}. LLMs have enabled the creation of more sophisticated AI models that can process complex data and provide accurate results. However, these models require significant computational resources, which can be a challenge in edge devices with limited resources \cite{ref166}. To address this issue, researchers have been exploring techniques such as model compression, pruning, and knowledge distillation to reduce the computational requirements of LLMs \cite{ref167}. These techniques are essential for enabling the deployment of LLMs in edge devices.

Edge AI has numerous applications across various industries, including healthcare, finance, and transportation \cite{ref168}. In healthcare, Edge AI can be used to analyze medical images and provide real-time diagnoses, enabling healthcare professionals to make timely decisions \cite{ref97}. In finance, Edge AI can be used to detect anomalies in transactions and prevent fraudulent activities \cite{ref169}. In transportation, Edge AI can be used to optimize traffic flow and reduce congestion \cite{ref170}. These applications demonstrate the potential of Edge AI to transform various industries and improve the efficiency of various processes.

Despite the numerous benefits of Edge AI, there are several challenges that need to be addressed. One of the significant challenges is the limited computational resources available in edge devices \cite{ref171}. Edge devices often have limited processing power, memory, and storage, which can make it challenging to deploy complex AI models. To address this issue, researchers have been exploring techniques such as edge-cloud collaboration, where edge devices can offload computationally intensive tasks to the cloud \cite{ref172}. Another challenge in Edge AI is the need for real-time processing and decision-making \cite{ref173}. Edge AI applications often require immediate processing and decision-making, which can be challenging in environments with limited computational resources. To address this issue, researchers have been exploring techniques such as real-time scheduling and resource allocation \cite{ref174}.

In conclusion, Edge AI is a revolutionary paradigm that has the potential to transform various industries by providing low-latency and real-time processing capabilities. The history of Edge AI dates back to the early 2000s, and it has gained significant traction in recent years due to the advent of edge computing and the proliferation of IoT devices. The importance of Edge AI lies in its ability to provide reduced latency and enable real-time decision-making, making it an essential component in applications such as autonomous vehicles, smart cities, and industrial automation. As Edge AI continues to evolve, it is likely to play a critical role in enabling the widespread adoption of AI in various industries, and addressing the challenges associated with its deployment will be crucial for realizing its full potential \cite{ref175}.

\subsection{Architectures and Techniques for Edge AI}

Architectures and techniques for Edge AI have been rapidly evolving to address the challenges of deploying artificial intelligence at the edge, as discussed in the previous section. One of the key techniques used in Edge AI is federated learning \cite{ref176,ref177}, which enables multiple devices to collaboratively train a model without sharing their data. This approach has been shown to be effective in reducing communication overhead and improving model accuracy \cite{ref116,ref178}, making it an essential component of Edge AI systems.

In addition to federated learning, other techniques such as distributed computation \cite{ref179,ref180} have been explored to improve the efficiency of Edge AI systems. Distributed computation involves distributing computational tasks across multiple devices to improve processing efficiency, and can be used in conjunction with federated learning to further enhance the performance of Edge AI systems. For example, \cite{ref181} proposes a knowledge cache-driven federated learning architecture that leverages distributed computation to improve the efficiency of model training.

To further optimize the performance of Edge AI systems, compression schemes \cite{ref182,ref183} have been widely adopted to reduce the size of models and improve communication efficiency. These schemes can be used to reduce the size of models, making them more suitable for deployment on resource-constrained edge devices. For instance, \cite{ref177} proposes a channel-wise model pruning algorithm that shrinks the footprint of local models while accounting for both data and system heterogeneity.

Conditional computation is another technique used in Edge AI to improve the efficiency of model inference \cite{ref184,ref185}. This approach involves dynamically adjusting the computational resources allocated to different tasks based on their priority and urgency, enabling Edge AI systems to optimize their performance in real-time. For example, \cite{ref186} proposes a federated intermediate layers learning approach that utilizes conditional computation to improve the efficiency of model training.

The development of Edge AI architectures has also been an active area of research, with various proposals aiming to provide a scalable and secure framework for deploying Edge AI applications. These include the use of blockchain \cite{ref116,ref187} and cloud-edge collaboration \cite{ref188,ref189}, which can provide a robust and efficient infrastructure for Edge AI systems.

The applications of Edge AI are diverse and widespread, and the techniques and architectures discussed in this section have been explored in various domains, including smart cities \cite{ref190,ref191}, healthcare \cite{ref192,ref193}, and industrial automation \cite{ref171,ref194}. As Edge AI continues to evolve, it is likely to play a critical role in enabling the widespread adoption of AI in these domains, and the techniques and architectures discussed in this section will be essential for realizing its full potential \cite{ref195,ref196}. The following section will delve into the applications of Edge AI, exploring its potential to transform various industries and aspects of our lives.

\subsection{Applications and Challenges of Edge AI}

The applications of Edge AI are diverse and widespread, transforming various industries and aspects of our lives. As discussed in the previous section, Edge AI has the potential to revolutionize numerous domains, including smart cities, healthcare, and industrial automation, by providing real-time processing and analysis of data. One of the primary applications of Edge AI is in the Internet of Things (IoT) \cite{ref3}, where it enables real-time processing and analysis of data generated by IoT devices. This is particularly useful in applications such as smart homes, industrial automation, and transportation systems, where timely decision-making is crucial. For instance, Edge AI can be used to detect anomalies in sensor data from industrial equipment, allowing for predictive maintenance and reducing downtime \cite{ref3}.

Another significant application of Edge AI is in smart cities \cite{ref25}, where it can be used to optimize traffic flow, energy consumption, and public safety. Edge AI can analyze data from traffic cameras, sensors, and other sources to detect accidents, traffic congestion, and other incidents, enabling swift response and mitigation. Additionally, Edge AI can be used to optimize energy consumption in smart buildings, reducing waste and improving efficiency \cite{ref25}. The use of Edge AI in smart cities is a natural extension of its application in IoT, as it enables the creation of more efficient, responsive, and sustainable urban environments.

In the healthcare sector, Edge AI \cite{ref197} can be used to analyze medical images, patient data, and other health-related information in real-time, enabling timely diagnosis and treatment. Edge AI can also be used to monitor patients remotely, reducing the need for hospitalization and improving patient outcomes. For example, Edge AI can be used to analyze data from wearable devices, such as heart rate and blood pressure monitors, to detect early warning signs of health problems \cite{ref197}. The application of Edge AI in healthcare is critical, as it has the potential to improve patient care, reduce costs, and enhance the overall quality of healthcare services.

Gaming is another area where Edge AI \cite{ref79} can have a significant impact, enabling real-time rendering, physics simulations, and AI-powered game mechanics. Edge AI can also be used to personalize gaming experiences, adapting to individual players' preferences and playing styles. For instance, Edge AI can be used to generate personalized game content, such as customized levels and characters, in real-time \cite{ref79}. The use of Edge AI in gaming is a testament to its versatility and potential to transform various industries and aspects of our lives.

However, despite the numerous benefits and applications of Edge AI, there are also several challenges that need to be addressed. One of the primary challenges is the need for specialized hardware and software to support Edge AI applications \cite{ref10}. This can be a significant barrier to adoption, particularly for small and medium-sized enterprises. Another challenge is the need for large amounts of high-quality training data to develop accurate Edge AI models \cite{ref198}. Security is also a significant concern in Edge AI applications \cite{ref44}, as sensitive data is often processed and stored at the edge. This requires robust security measures, such as encryption, access control, and intrusion detection, to protect against cyber threats.

Furthermore, Edge AI applications often require significant computational resources, which can be a challenge in resource-constrained edge environments \cite{ref199}. This requires careful optimization of Edge AI models and algorithms to minimize computational requirements while maintaining accuracy and performance. Energy efficiency is also a critical concern in Edge AI applications \cite{ref166}, as many edge devices are battery-powered or have limited power budgets. By addressing these challenges and developing innovative Edge AI solutions, we can unlock the full potential of Edge AI and create a more efficient, responsive, and intelligent world \cite{ref78}. As we move forward, it is essential to continue exploring the applications, challenges, and opportunities of Edge AI, and to develop new technologies and strategies that can help overcome the challenges and realize the full potential of Edge AI.

\section{Challenges, Open Research Directions, and Future Outlook}

\subsection{Introduction to Challenges in Edge Computing}

Edge computing has emerged as a promising technology to reduce latency and improve real-time processing capabilities by bringing computing resources closer to the end-users. However, despite its potential, edge computing faces several challenges that need to be addressed to ensure its widespread adoption. In this subsection, we will introduce the current challenges in edge computing, including scalability, reliability, and standardization, and provide an overview of the importance of addressing these challenges.

One of the significant challenges in edge computing is scalability \cite{ref2}. As the number of edge devices and applications increases, the edge infrastructure needs to be able to scale up to handle the growing demand. However, this can be a complex task, especially when considering the heterogeneous nature of edge devices and the varying requirements of different applications. To address this challenge, researchers have proposed various solutions, such as distributed edge computing architectures \cite{ref200} and edge-based machine learning algorithms \cite{ref43}. These solutions aim to improve the scalability of edge computing by enabling the efficient distribution of resources and tasks across multiple edge devices.

Another challenge in edge computing is reliability \cite{ref18}. Edge devices are often resource-constrained and may be prone to failures, which can lead to service disruptions and affect the overall quality of experience. To ensure reliable edge computing, researchers have proposed various solutions, such as fault-tolerant edge computing architectures \cite{ref201} and edge-based fault detection and recovery mechanisms \cite{ref150}. These solutions aim to improve the reliability of edge computing by detecting and recovering from failures in a timely and efficient manner.

Standardization is another significant challenge in edge computing \cite{ref41}. The lack of standardization in edge computing can lead to interoperability issues and make it difficult to develop and deploy edge applications. To address this challenge, researchers have proposed various solutions, such as standardized edge computing architectures \cite{ref202} and edge-based APIs \cite{ref15}. These solutions aim to improve the standardization of edge computing by providing a common framework for developing and deploying edge applications.

In addition to these challenges, edge computing also faces several other challenges, such as security \cite{ref16}, privacy \cite{ref107}, and energy efficiency \cite{ref203}. These challenges are critical to the widespread adoption of edge computing and need to be addressed through innovative solutions and technologies.

The importance of addressing these challenges cannot be overstated, as edge computing has the potential to revolutionize various industries, such as healthcare, transportation, and manufacturing, by providing real-time processing and analysis of data. If the challenges in edge computing are not addressed, it can lead to significant consequences, such as service disruptions, security breaches, and energy inefficiencies. Therefore, it is essential to invest in research and development to address these challenges and ensure the widespread adoption of edge computing.

To overcome these challenges, researchers have made significant progress in recent years. For example, researchers have proposed various solutions to improve the scalability of edge computing, such as distributed edge computing architectures \cite{ref204} and edge-based machine learning algorithms \cite{ref7}. Additionally, researchers have proposed various solutions to improve the reliability of edge computing, such as fault-tolerant edge computing architectures \cite{ref205} and edge-based fault detection and recovery mechanisms \cite{ref206}. These advancements have paved the way for the integration of edge computing with emerging technologies, such as 6G, artificial intelligence, and the Internet of Things (IoT), which will be discussed in the following section.

In conclusion, edge computing faces several challenges, including scalability, reliability, and standardization, that need to be addressed to ensure its widespread adoption. These challenges are critical to the success of edge computing and need to be addressed through innovative solutions and technologies. By addressing these challenges, we can unlock the full potential of edge computing and revolutionize various industries by providing real-time processing and analysis of data, ultimately leading to the development of new applications and services that will be enabled by the integration of edge computing with emerging technologies.

\subsection{Open Research Directions in Edge Computing}

The field of edge computing is rapidly evolving, and several open research directions are emerging as the technology continues to advance. Building on the challenges and limitations discussed in the previous section, researchers are now focusing on the integration of edge computing with emerging technologies such as 6G, artificial intelligence, and the Internet of Things (IoT). As noted in \cite{ref207}, the convergence of edge computing and 6G networks is expected to enable a wide range of new applications and services, including immersive media, smart cities, and industrial automation.

One of the key areas of focus is the integration of edge computing with artificial intelligence, as highlighted in \cite{ref28}. The use of AI and machine learning algorithms at the edge of the network can enable real-time processing and analysis of data, reducing latency and improving overall system performance. For example, \cite{ref7} discusses the use of deep learning algorithms at the edge of the network, enabling applications such as real-time object detection and classification. This integration is crucial in addressing the challenges of scalability, reliability, and standardization discussed earlier, as it enables the efficient distribution of resources and tasks across multiple edge devices.

Another important area of research is the integration of edge computing with the IoT, as noted in \cite{ref24}. The IoT is expected to generate a vast amount of data, which will need to be processed and analyzed in real-time. Edge computing can play a key role in enabling this, by providing a distributed computing infrastructure that can process data closer to the source. For example, \cite{ref3} discusses the use of edge computing in IoT applications such as smart cities and industrial automation. This integration has the potential to transform various industries, including healthcare, transportation, and manufacturing, by providing real-time processing and analysis of data.

The use of edge computing in 6G networks is also a key area of research, as highlighted in \cite{ref208}. The use of edge computing in 6G networks can enable a wide range of new applications and services, including immersive media, smart cities, and industrial automation. For example, \cite{ref209} discusses the use of mobile edge computing in 6G networks, enabling applications such as virtual and augmented reality. This convergence of edge computing and 6G networks is expected to have a significant impact on various industries and society as a whole, as discussed in the following section.

In addition to these areas, there are several other open research directions in edge computing, including the development of new edge computing architectures, the use of edge computing in vehicular networks, and the integration of edge computing with other emerging technologies such as blockchain and quantum computing. For example, \cite{ref210} discusses the use of a three-tier architecture for cognitive edge computing. The development of new edge computing architectures is crucial in addressing the challenges of scalability, reliability, and standardization, as it enables the efficient distribution of resources and tasks across multiple edge devices.

The use of edge computing in vehicular networks is also a key area of research, as highlighted in \cite{ref211}. The use of edge computing in vehicular networks can enable a wide range of new applications and services, including smart transportation systems and autonomous vehicles. For example, \cite{ref212} discusses the use of mobility-enhanced edge intelligence for smart and green 6G networks. This integration has the potential to transform the transportation industry by providing real-time processing and analysis of data.

The integration of edge computing with blockchain is also a key area of research, as highlighted in \cite{ref209}. The use of blockchain in edge computing can enable secure and transparent data processing and analysis, reducing the risk of data breaches and cyber attacks. For example, \cite{ref213} discusses the use of blockchain in fog-IoT architecture, enabling secure and transparent data processing and analysis. This integration is crucial in addressing the challenges of security and privacy discussed earlier, as it enables the secure and transparent processing and analysis of data.

In conclusion, the field of edge computing is rapidly evolving, and several open research directions are emerging as the technology continues to advance. The integration of edge computing with emerging technologies such as 6G, artificial intelligence, and the IoT is expected to enable a wide range of new applications and services, including immersive media, smart cities, and industrial automation. The development of new edge computing architectures, the use of edge computing in vehicular networks, and the integration of edge computing with other emerging technologies such as blockchain and quantum computing are also key areas of research. As noted in \cite{ref214}, the use of edge computing can enable real-time processing and analysis of data, reducing latency and improving overall system performance, and its potential benefits will be discussed in the following section.

\subsection{Future Outlook and Potential Impact of Edge Computing}

The future outlook of edge computing is promising, with potential impacts on various industries and society as a whole. As edge computing continues to evolve, it is expected to play a crucial role in enabling new use cases and applications that require low latency, high bandwidth, and real-time processing \cite{ref1}. Building on the open research directions and emerging technologies discussed earlier, such as the integration of edge computing with 6G, artificial intelligence, and the IoT, the widespread adoption of edge computing is expected to transform industries such as healthcare, manufacturing, transportation, and education, among others.

In healthcare, for instance, edge computing can enable real-time analysis of medical images, remote patient monitoring, and personalized medicine \cite{ref215}. By analyzing medical images in real-time, edge computing can enable faster diagnosis and treatment, while remote patient monitoring can allow healthcare professionals to track patient vital signs and respond quickly to emergencies. Similarly, in manufacturing, edge computing can enable predictive maintenance, quality control, and supply chain optimization \cite{ref159}, by analyzing sensor data from machines in real-time and reducing downtime.

The potential applications of edge computing are vast, and its impact can be seen in various industries, including transportation, where it can enable autonomous vehicles, smart traffic management, and route optimization \cite{ref19}, and education, where it can enable personalized learning, virtual classrooms, and remote education \cite{ref216}. However, the widespread adoption of edge computing also poses several challenges, including security, privacy, and scalability \cite{ref16}, which need to be addressed through the development of new technologies and architectures.

To address these challenges, researchers and industry experts are exploring new technologies and architectures, such as federated learning, distributed computing, and blockchain \cite{ref217}. These technologies can enable secure and transparent data sharing and management, while also providing real-time processing and analysis capabilities. As edge computing continues to evolve, it is expected to play a crucial role in enabling new use cases and applications, and its potential benefits include improved efficiency, productivity, and decision-making, as well as new business models and revenue streams.

In conclusion, the future outlook of edge computing is promising, with potential impacts on various industries and society as a whole. However, the widespread adoption of edge computing also poses several challenges, including security, privacy, and scalability \cite{ref218}. To address these challenges, researchers and industry experts are exploring new technologies and architectures, and as edge computing continues to evolve, it is expected to play a crucial role in enabling new use cases and applications that require low latency, high bandwidth, and real-time processing \cite{ref24}. Ultimately, the development of a comprehensive roadmap for sustainable edge computing, as discussed in the next section, will be crucial in ensuring that the benefits of edge computing are realized in a responsible and sustainable manner.

\subsection{Roadmap for Sustainable Edge Computing Development}

To achieve a sustainable and widely adopted edge computing ecosystem, a comprehensive roadmap is essential, building on the potential impacts and challenges of edge computing discussed earlier. This roadmap should outline the key steps and milestones necessary for the development and implementation of sustainable edge computing, addressing the challenges and open issues in edge computing architecture, such as scalability, reliability, and security, which are critical for ensuring the widespread adoption of edge computing \cite{ref1}.

As edge computing continues to evolve and play a crucial role in enabling new use cases and applications, the roadmap should prioritize the development of sustainable edge computing architectures, leveraging emerging technologies such as fog computing, mobile edge computing, and edge-based machine learning to overcome the challenges of scalability, reliability, and security. Furthermore, the roadmap should focus on resource management and optimization techniques, including resource allocation, scheduling, and energy efficiency, which are vital for minimizing latency and reducing energy consumption, and can be achieved through the use of machine learning and other optimization techniques.

In addition to these technical aspects, the roadmap should also consider the societal and environmental implications of edge computing, as discussed in \cite{ref203}, including the potential impact of edge computing on energy consumption, e-waste, and carbon emissions, as well as its potential to support sustainable development and reduce the environmental footprint of computing systems. The roadmap should prioritize security and privacy in edge computing, including data protection, authentication, and access control, which are critical for ensuring the trustworthiness and reliability of edge computing systems, and can be achieved through the use of encryption, authentication, and access control mechanisms.

The development of edge AI is also critical for enabling the widespread adoption of edge computing, as emphasized in \cite{ref195}, and should be included in the roadmap, including the development of novel AI methods and capabilities, as well as the integration of AI with edge computing architectures. By outlining a series of milestones and key steps, the roadmap can help ensure that edge computing is developed and implemented in a sustainable and responsible manner, with a focus on minimizing its environmental footprint and supporting sustainable development, ultimately leading to improved efficiency, productivity, and decision-making, as well as new business models and revenue streams.

In conclusion, the roadmap for sustainable edge computing development should prioritize the development of sustainable edge computing architectures, resource management and optimization techniques, security and privacy mechanisms, and edge AI, and should be developed through collaboration and knowledge sharing among researchers, developers, and industry stakeholders, as discussed in \cite{ref219}, to ensure that the benefits of edge computing are realized in a responsible and sustainable manner, and to support the future of computing and networking research.


\begin{thebibliography}{219}
\addcontentsline{toc}{section}{References}
\bibitem{ref1} Edge Computing: A Comprehensive Survey of Current Initiatives and a Roadmap for a Sustainable Edge Computing Development. arXiv:1912.08530, 2019. \url{https://arxiv.org/abs/1912.08530}
\bibitem{ref2} Edge Computing: A Systematic Mapping Study. arXiv:2102.02720, 2021. \url{https://arxiv.org/abs/2102.02720}
\bibitem{ref3} Edge Computing for IoT. arXiv:2402.13056, 2024. \url{https://arxiv.org/abs/2402.13056}
\bibitem{ref4} Edge Intelligence: Architectures, Challenges, and Applications. arXiv:2003.12172, 2020. \url{https://arxiv.org/abs/2003.12172}
\bibitem{ref5} A Comprehensive Review and a Taxonomy of Edge Machine Learning: Requirements, Paradigms, and Techniques. arXiv:2302.08571, 2023. \url{https://arxiv.org/abs/2302.08571}
\bibitem{ref6} Deep Learning in the Era of Edge Computing: Challenges and Opportunities. arXiv:2010.08861, 2020. \url{https://arxiv.org/abs/2010.08861}
\bibitem{ref7} Deep Learning at the Edge. arXiv:1910.10231, 2019. \url{https://arxiv.org/abs/1910.10231}
\bibitem{ref8} Mobile Edge Computing. arXiv:2404.17944, 2024. \url{https://arxiv.org/abs/2404.17944}
\bibitem{ref9} Edge computing vs centralized cloud: Impact of communication latency on the energy consumption of LTE terminal nodes. arXiv:2111.10076, 2021. \url{https://arxiv.org/abs/2111.10076}
\bibitem{ref10} AI on the Edge: Rethinking AI-based IoT Applications Using Specialized Edge Architectures. arXiv:2003.12488, 2020. \url{https://arxiv.org/abs/2003.12488}
\bibitem{ref11} The Benefits of Edge Computing in Healthcare, Smart Cities, and IoT. arXiv:2112.01250, 2021. \url{https://arxiv.org/abs/2112.01250}
\bibitem{ref12} At the Edge of a Seamless Cloud Experience. arXiv:2111.06157, 2021. \url{https://arxiv.org/abs/2111.06157}
\bibitem{ref13} Sustainable AI Processing at the Edge. arXiv:2207.01209, 2022. \url{https://arxiv.org/abs/2207.01209}
\bibitem{ref14} An approach to define Very High Capacity Networks with improved quality at an affordable cost. arXiv:2011.03685, 2020. \url{https://arxiv.org/abs/2011.03685}
\bibitem{ref15} EaaS: A Service-Oriented Edge Computing Framework Towards Distributed Intelligence. arXiv:2209.06613, 2022. \url{https://arxiv.org/abs/2209.06613}
\bibitem{ref16} Edge Security: Challenges and Issues. arXiv:2206.07164, 2022. \url{https://arxiv.org/abs/2206.07164}
\bibitem{ref17} Review of Data Integrity Attacks and Mitigation Methods in Edge Computing. arXiv:2112.02757, 2021. \url{https://arxiv.org/abs/2112.02757}
\bibitem{ref18} Dependability in Edge Computing. arXiv:1710.11222, 2017. \url{https://arxiv.org/abs/1710.11222}
\bibitem{ref19} Wireless Edge Computing With Latency and Reliability Guarantees. arXiv:1905.05316, 2019. \url{https://arxiv.org/abs/1905.05316}
\bibitem{ref20} Dynamic DAG-Application Scheduling for Multi-Tier Edge Computing in Heterogeneous Networks. arXiv:2409.10839, 2024. \url{https://arxiv.org/abs/2409.10839}
\bibitem{ref21} Edge computing for cyber-physical systems: A systematic mapping study emphasizing trustworthiness. arXiv:2112.00619, 2021. \url{https://arxiv.org/abs/2112.00619}
\bibitem{ref22} Towards an Intelligent Edge: Wireless Communication Meets Machine Learning. arXiv:1809.00343, 2018. \url{https://arxiv.org/abs/1809.00343}
\bibitem{ref23} Mobile Edge Computing-Enabled Heterogeneous Networks. arXiv:1804.07756, 2018. \url{https://arxiv.org/abs/1804.07756}
\bibitem{ref24} Edge Computing in IoT: A 6G Perspective. arXiv:2111.08943, 2021. \url{https://arxiv.org/abs/2111.08943}
\bibitem{ref25} Edge-Computing-Enabled Smart Cities: A Comprehensive Survey. arXiv:1909.08747, 2019. \url{https://arxiv.org/abs/1909.08747}
\bibitem{ref26} A Survey on UAV-enabled Edge Computing: Resource Management Perspective. arXiv:2210.06679, 2022. \url{https://arxiv.org/abs/2210.06679}
\bibitem{ref27} A Survey on Open-Source Edge Computing Simulators and Emulators: The Computing and Networking Convergence Perspective. arXiv:2505.09995, 2025. \url{https://arxiv.org/abs/2505.09995}
\bibitem{ref28} Edge AI: A Taxonomy, Systematic Review and Future Directions. arXiv:2407.04053, 2024. \url{https://arxiv.org/abs/2407.04053}
\bibitem{ref29} What is a Fog Node A Tutorial on Current Concepts towards a Common Definition. arXiv:1611.09193, 2016. \url{https://arxiv.org/abs/1611.09193}
\bibitem{ref30} EdgeFlow: Open-Source Multi-layer Data Flow Processing in Edge Computing for 5G and Beyond. arXiv:1801.02206, 2018. \url{https://arxiv.org/abs/1801.02206}
\bibitem{ref31} DHT-based Edge and Fog Computing Systems: Infrastructures and Applications. arXiv:2211.05851, 2022. \url{https://arxiv.org/abs/2211.05851}
\bibitem{ref32} V-Edge: Virtual Edge Computing as an Enabler for Novel Microservices and Cooperative Computing. arXiv:2106.10063, 2021. \url{https://arxiv.org/abs/2106.10063}
\bibitem{ref33} Resource Management in Edge and Fog Computing using FogBus2 Framework. arXiv:2108.00591, 2021. \url{https://arxiv.org/abs/2108.00591}
\bibitem{ref34} Fog Function: Serverless Fog Computing for Data Intensive IoT Services. arXiv:1907.08278, 2019. \url{https://arxiv.org/abs/1907.08278}
\bibitem{ref35} Edge-PRUNE: Flexible Distributed Deep Learning Inference. arXiv:2204.12947, 2022. \url{https://arxiv.org/abs/2204.12947}
\bibitem{ref36} Edge-Cloud Collaborative Computing on Distributed Intelligence and Model Optimization: A Survey. arXiv:2505.01821, 2025. \url{https://arxiv.org/abs/2505.01821}
\bibitem{ref37} Security Assessment and Hardening of Fog Computing Systems. arXiv:2308.12707, 2023. \url{https://arxiv.org/abs/2308.12707}
\bibitem{ref38} Trustworthy Edge Computing through Blockchains. arXiv:2005.07741, 2020. \url{https://arxiv.org/abs/2005.07741}
\bibitem{ref39} How BlockChain Can Help Enhance The Security And Privacy in Edge Computing?. arXiv:2111.00416, 2021. \url{https://arxiv.org/abs/2111.00416}
\bibitem{ref40} The Security and Privacy of Mobile Edge Computing: An Artificial Intelligence Perspective. arXiv:2401.01589, 2024. \url{https://arxiv.org/abs/2401.01589}
\bibitem{ref41} A Taxonomy for Management and Optimization of Multiple Resources in Edge Computing. arXiv:1801.05610, 2018. \url{https://arxiv.org/abs/1801.05610}
\bibitem{ref42} Proactive and Reactive Autoscaling Techniques for Edge Computing. arXiv:2510.10166, 2025. \url{https://arxiv.org/abs/2510.10166}
\bibitem{ref43} EdgeAI: A Vision for Deep Learning in IoT Era. arXiv:1910.10356, 2019. \url{https://arxiv.org/abs/1910.10356}
\bibitem{ref44} The AI Shadow War: SaaS vs. Edge Computing Architectures. arXiv:2507.11545, 2025. \url{https://arxiv.org/abs/2507.11545}
\bibitem{ref45} Using cloud and fog computing for large scale IoT-based urban sound classification. arXiv:1910.07652, 2019. \url{https://arxiv.org/abs/1910.07652}
\bibitem{ref46} Deep Learning for Reliable Mobile Edge Analytics in Intelligent Transportation Systems. arXiv:1712.04135, 2017. \url{https://arxiv.org/abs/1712.04135}
\bibitem{ref47} Intelligent Energy Management with IoT Framework in Smart Cities Using Intelligent Analysis: An Application of Machine Learning Methods for Complex Networks and Systems. arXiv:2306.05567, 2023. \url{https://arxiv.org/abs/2306.05567}
\bibitem{ref48} Fog Computing for Sustainable Smart Cities: A Survey. arXiv:1703.07079, 2017. \url{https://arxiv.org/abs/1703.07079}
\bibitem{ref49} Green Internet of Vehicles (IoV) in the 6G Era: Toward Sustainable Vehicular Communications and Networking. arXiv:2108.11879, 2021. \url{https://arxiv.org/abs/2108.11879}
\bibitem{ref50} The Role of LLMs in Sustainable Smart Cities: Applications, Challenges, and Future Directions. arXiv:2402.14596, 2024. \url{https://arxiv.org/abs/2402.14596}
\bibitem{ref51} Role of Edge Device and Cloud Machine Learning in Point-of-Care Solutions Using Imaging Diagnostics for Population Screening. arXiv:2006.13808, 2020. \url{https://arxiv.org/abs/2006.13808}
\bibitem{ref52} Ubiquitous Mobile Health Monitoring System for Elderly (UMHMSE). arXiv:1107.3695, 2011. \url{https://arxiv.org/abs/1107.3695}
\bibitem{ref53} Optimizing Edge Gaming Slices through an Enhanced User Plane Function and Analytics in Beyond-5G Networks. arXiv:2507.17843, 2025. \url{https://arxiv.org/abs/2507.17843}
\bibitem{ref54} Poster: Enabling Flexible Edge-assisted XR. arXiv:2309.04548, 2023. \url{https://arxiv.org/abs/2309.04548}
\bibitem{ref55} Artificial Intelligence at the Edge. arXiv:2012.05410, 2020. \url{https://arxiv.org/abs/2012.05410}
\bibitem{ref56} IoT-Enabled Low-Cost Fog Computing System with Online Machine Learning for Accurate and Low-Latency Heart Monitoring in Rural Healthcare Settings. arXiv:2302.14131, 2023. \url{https://arxiv.org/abs/2302.14131}
\bibitem{ref57} A Survey on Security and Privacy Issues in Modern Healthcare Systems: Attacks and Defenses. arXiv:2005.07359, 2020. \url{https://arxiv.org/abs/2005.07359}
\bibitem{ref58} PRISM: Privacy Preserving Healthcare Internet of Things Security Management. arXiv:2212.14736, 2022. \url{https://arxiv.org/abs/2212.14736}
\bibitem{ref59} Edge Computing for Smart Health: Context-Aware Approaches, Opportunities, and Challenges. arXiv:2004.07311, 2020. \url{https://arxiv.org/abs/2004.07311}
\bibitem{ref60} EdgeLaaS: Edge Learning as a Service for Knowledge-Centric Connected Healthcare. arXiv:1808.09141, 2018. \url{https://arxiv.org/abs/1808.09141}
\bibitem{ref61} Autonomy and Intelligence in the Computing Continuum: Challenges, Enablers, and Future Directions for Orchestration. arXiv:2205.01423, 2022. \url{https://arxiv.org/abs/2205.01423}
\bibitem{ref62} Software Defined Edge Computing for Distributed Management and Scalable Control in IoT Multinetworks. arXiv:2104.02426, 2021. \url{https://arxiv.org/abs/2104.02426}
\bibitem{ref63} Internet of Things (IoT) and New Computing Paradigms. arXiv:1812.00591, 2018. \url{https://arxiv.org/abs/1812.00591}
\bibitem{ref64} A Survey on Edge Benchmarking. arXiv:2004.11725, 2020. \url{https://arxiv.org/abs/2004.11725}
\bibitem{ref65} Towards Enabling Novel Edge-Enabled Applications. arXiv:1805.06989, 2018. \url{https://arxiv.org/abs/1805.06989}
\bibitem{ref66} Joint Computing Offloading and Resource Allocation for Classification Intelligent Tasks in MEC Systems. arXiv:2307.03205, 2023. \url{https://arxiv.org/abs/2307.03205}
\bibitem{ref67} Joint Task Offloading Scheduling and Transmit Power Allocation for Mobile-Edge Computing Systems. arXiv:1701.05055, 2017. \url{https://arxiv.org/abs/1701.05055}
\bibitem{ref68} TPTO: A Transformer-PPO based Task Offloading Solution for Edge Computing Environments. arXiv:2312.11739, 2023. \url{https://arxiv.org/abs/2312.11739}
\bibitem{ref69} Multiuser Resource Allocation for Mobile-Edge Computation Offloading. arXiv:1604.02519, 2016. \url{https://arxiv.org/abs/1604.02519}
\bibitem{ref70} Joint Optimization of Execution Latency and Energy Consumption for Mobile Edge Computing with Data Compression and Task Allocation. arXiv:1909.12695, 2019. \url{https://arxiv.org/abs/1909.12695}
\bibitem{ref71} Federated Learning Assisted Deep Q-Learning for Joint Task Offloading and Fronthaul Segment Routing in Open RAN. arXiv:2308.07258, 2023. \url{https://arxiv.org/abs/2308.07258}
\bibitem{ref72} Latency and Reliability-Aware Task Offloading and Resource Allocation for Mobile Edge Computing. arXiv:1710.00590, 2017. \url{https://arxiv.org/abs/1710.00590}
\bibitem{ref73} Energy-Efficient Deadline-Aware Edge Computing: Bandit Learning with Partial Observations in Multi-Channel Systems. arXiv:2308.06647, 2023. \url{https://arxiv.org/abs/2308.06647}
\bibitem{ref74} A Novel Cross Entropy Approach for Offloading Learning in Mobile Edge Computing. arXiv:1912.03838, 2019. \url{https://arxiv.org/abs/1912.03838}
\bibitem{ref75} EdgeConvEns: Convolutional Ensemble Learning for Edge Intelligence. arXiv:2307.14381, 2023. \url{https://arxiv.org/abs/2307.14381}
\bibitem{ref76} Federated Deep Reinforcement Learning. arXiv:1901.08277, 2019. \url{https://arxiv.org/abs/1901.08277}
\bibitem{ref77} Coded Computing for Federated Learning at the Edge. arXiv:2007.03273, 2020. \url{https://arxiv.org/abs/2007.03273}
\bibitem{ref78} Edge Intelligence: Paving the Last Mile of Artificial Intelligence with Edge Computing. arXiv:1905.10083, 2019. \url{https://arxiv.org/abs/1905.10083}
\bibitem{ref79} Analyzing Transformer Models and Knowledge Distillation Approaches for Image Captioning on Edge AI. arXiv:2506.03607, 2025. \url{https://arxiv.org/abs/2506.03607}
\bibitem{ref80} In-Edge AI: Intelligentizing Mobile Edge Computing, Caching and Communication by Federated Learning. arXiv:1809.07857, 2018. \url{https://arxiv.org/abs/1809.07857}
\bibitem{ref81} Conic Optimization Theory: Convexification Techniques and Numerical Algorithms. arXiv:1709.08841, 2017. \url{https://arxiv.org/abs/1709.08841}
\bibitem{ref82} Resource Dimensioning for Single-Cell Edge Video Analytics. arXiv:2305.05568, 2023. \url{https://arxiv.org/abs/2305.05568}
\bibitem{ref83} Edge Delayed Deep Deterministic Policy Gradient: efficient continuous control for edge scenarios. arXiv:2412.06390, 2024. \url{https://arxiv.org/abs/2412.06390}
\bibitem{ref84} Dynamic Offloading Loading Optimization in distributed Fault Diagnosis system with Deep Reinforcement Learning Approach. arXiv:2103.02174, 2021. \url{https://arxiv.org/abs/2103.02174}
\bibitem{ref85} Social Welfare Maximization for Collaborative Edge Computing: A Deep Reinforcement Learning-Based Approach. arXiv:2211.06861, 2022. \url{https://arxiv.org/abs/2211.06861}
\bibitem{ref86} POSEIDON : Efficient Function Placement at the Edge using Deep Reinforcement Learning. arXiv:2410.11879, 2024. \url{https://arxiv.org/abs/2410.11879}
\bibitem{ref87} Fast Convex Optimization for Two-Layer ReLU Networks: Equivalent Model Classes and Cone Decompositions. arXiv:2202.01331, 2022. \url{https://arxiv.org/abs/2202.01331}
\bibitem{ref88} Tutorials on Advanced Optimization Methods. arXiv:2007.13545, 2020. \url{https://arxiv.org/abs/2007.13545}
\bibitem{ref89} Edge Computing in Transportation: Security Issues and Challenges. arXiv:2012.11206, 2020. \url{https://arxiv.org/abs/2012.11206}
\bibitem{ref90} EdgeLeakage: Membership Information Leakage in Distributed Edge Intelligence Systems. arXiv:2404.16851, 2024. \url{https://arxiv.org/abs/2404.16851}
\bibitem{ref91} Physical and Software Based Fault Injection Attacks Against TEEs in Mobile Devices: A Systemisation of Knowledge. arXiv:2411.14878, 2024. \url{https://arxiv.org/abs/2411.14878}
\bibitem{ref92} Security in Mobile Edge Caching with Reinforcement Learning. arXiv:1801.05915, 2018. \url{https://arxiv.org/abs/1801.05915}
\bibitem{ref93} A Novel Access Control and Privacy-Enhancing Approach for Models in Edge Computing. arXiv:2411.03847, 2024. \url{https://arxiv.org/abs/2411.03847}
\bibitem{ref94} Private Edge Computing for Linear Inference Based on Secret Sharing. arXiv:2005.08593, 2020. \url{https://arxiv.org/abs/2005.08593}
\bibitem{ref95} A Lightweight and Privacy-Preserving Authentication Protocol for Mobile Edge Computing. arXiv:1907.08896, 2019. \url{https://arxiv.org/abs/1907.08896}
\bibitem{ref96} AEAKA: An Adaptive and Efficient Authentication and Key Agreement Scheme for IoT in Cloud-Edge-Device Collaborative Environments. arXiv:2411.09231, 2024. \url{https://arxiv.org/abs/2411.09231}
\bibitem{ref97} Blockchain Powered Edge Intelligence for U-Healthcare in Privacy Critical and Time Sensitive Environment. arXiv:2506.02038, 2025. \url{https://arxiv.org/abs/2506.02038}
\bibitem{ref98} EC-SVC: Secure CAN Bus In-Vehicle Communications with Fine-grained Access Control Based on Edge Computing. arXiv:2010.14747, 2020. \url{https://arxiv.org/abs/2010.14747}
\bibitem{ref99} Secure Coded Cooperative Computation at the Heterogeneous Edge against Byzantine Attacks. arXiv:1908.05385, 2019. \url{https://arxiv.org/abs/1908.05385}
\bibitem{ref100} Reconfigurable Security: Edge Computing-based Framework for IoT. arXiv:1709.06223, 2017. \url{https://arxiv.org/abs/1709.06223}
\bibitem{ref101} SPX: Preserving End-to-End Security for Edge Computing. arXiv:1809.09038, 2018. \url{https://arxiv.org/abs/1809.09038}
\bibitem{ref102} Privacy-Preserving Edge Caching: A Probabilistic Approach. arXiv:2208.00297, 2022. \url{https://arxiv.org/abs/2208.00297}
\bibitem{ref103} Smart Contract-based Computing ResourcesTrading in Edge Computing. arXiv:2006.15824, 2020. \url{https://arxiv.org/abs/2006.15824}
\bibitem{ref104} Joint Storage Allocation and Computation Design for Private Edge Computing. arXiv:2010.09012, 2020. \url{https://arxiv.org/abs/2010.09012}
\bibitem{ref105} Privacy-Preserving Edge Computing from Pairing-Based Inner Product Functional Encryption. arXiv:2504.02068, 2025. \url{https://arxiv.org/abs/2504.02068}
\bibitem{ref106} A lightweight blockchain-based access control scheme for integrated edge computing in the internet of things. arXiv:2111.06544, 2021. \url{https://arxiv.org/abs/2111.06544}
\bibitem{ref107} Edge Computing for IoT: Novel Insights from a Comparative Analysis of Access Control Models. arXiv:2405.07685, 2024. \url{https://arxiv.org/abs/2405.07685}
\bibitem{ref108} MQT-TZ: Hardening IoT Brokers Using ARM TrustZone. arXiv:2007.12442, 2020. \url{https://arxiv.org/abs/2007.12442}
\bibitem{ref109} Collaborative Inference over Wireless Channels with Feature Differential Privacy. arXiv:2410.19917, 2024. \url{https://arxiv.org/abs/2410.19917}
\bibitem{ref110} PA-iMFL: Communication-Efficient Privacy Amplification Method against Data Reconstruction Attack in Improved Multi-Layer Federated Learning. arXiv:2309.13864, 2023. \url{https://arxiv.org/abs/2309.13864}
\bibitem{ref111} Envisioning the Future of Cyber Security in Post-Quantum Era: A Survey on PQ Standardization, Applications, Challenges and Opportunities. arXiv:2310.12037, 2023. \url{https://arxiv.org/abs/2310.12037}
\bibitem{ref112} A Comprehensive Framework for Building Highly Secure, Network-Connected Devices: Chip to App. arXiv:2501.13716, 2025. \url{https://arxiv.org/abs/2501.13716}
\bibitem{ref113} Blockchain for Mobile Edge Computing: Consensus Mechanisms and Scalability. arXiv:2006.07578, 2020. \url{https://arxiv.org/abs/2006.07578}
\bibitem{ref114} Privacy Checklist: Privacy Violation Detection Grounding on Contextual Integrity Theory. arXiv:2408.10053, 2024. \url{https://arxiv.org/abs/2408.10053}
\bibitem{ref115} Vicious Classifiers: Assessing Inference-time Data Reconstruction Risk in Edge Computing. arXiv:2212.04223, 2022. \url{https://arxiv.org/abs/2212.04223}
\bibitem{ref116} Federated Learning Meets Blockchain in Edge Computing: Opportunities and Challenges. arXiv:2104.01776, 2021. \url{https://arxiv.org/abs/2104.01776}
\bibitem{ref117} Resource Scheduling in Edge Computing: A Survey. arXiv:2108.08059, 2021. \url{https://arxiv.org/abs/2108.08059}
\bibitem{ref118} Predictive Edge Computing with Hard Deadlines. arXiv:1805.12272, 2018. \url{https://arxiv.org/abs/1805.12272}
\bibitem{ref119} EASE: Energy-Aware Job Scheduling for Vehicular Edge Networks With Renewable Energy Resources. arXiv:2111.02186, 2021. \url{https://arxiv.org/abs/2111.02186}
\bibitem{ref120} Energy-Efficient Real-Time Job Mapping and Resource Management in Mobile-Edge Computing. arXiv:2506.12686, 2025. \url{https://arxiv.org/abs/2506.12686}
\bibitem{ref121} Joint User Scheduling and Computing Resource Allocation Optimization in Asynchronous Mobile Edge Computing Networks. arXiv:2401.11377, 2024. \url{https://arxiv.org/abs/2401.11377}
\bibitem{ref122} Fast and Adaptive Task Management in MEC: A Deep Learning Approach Using Pointer Networks. arXiv:2507.09346, 2025. \url{https://arxiv.org/abs/2507.09346}
\bibitem{ref123} AMP4EC: Adaptive Model Partitioning Framework for Efficient Deep Learning Inference in Edge Computing Environments. arXiv:2504.00407, 2025. \url{https://arxiv.org/abs/2504.00407}
\bibitem{ref124} EdgeTimer: Adaptive Multi-Timescale Scheduling in Mobile Edge Computing with Deep Reinforcement Learning. arXiv:2406.07342, 2024. \url{https://arxiv.org/abs/2406.07342}
\bibitem{ref125} Online Learning for Offloading and Autoscaling in Renewable-Powered Mobile Edge Computing. arXiv:1609.05087, 2016. \url{https://arxiv.org/abs/1609.05087}
\bibitem{ref126} Task Offloading Optimization in NOMA-Enabled Multi-hop Mobile Edge Computing System Using Conflict Graph. arXiv:2104.11801, 2021. \url{https://arxiv.org/abs/2104.11801}
\bibitem{ref127} Multi-objective Deep Reinforcement Learning for Mobile Edge Computing. arXiv:2307.14346, 2023. \url{https://arxiv.org/abs/2307.14346}
\bibitem{ref128} Joint Task Offloading and Resource Allocation for Multi-Server Mobile-Edge Computing Networks. arXiv:1705.00704, 2017. \url{https://arxiv.org/abs/1705.00704}
\bibitem{ref129} Computation Resource Allocation for Heterogeneous Time-Critical IoT Services in MEC. arXiv:2002.04851, 2020. \url{https://arxiv.org/abs/2002.04851}
\bibitem{ref130} A Hybrid Artificial Neural Network for Task Offloading in Mobile Edge Computing. arXiv:2206.06301, 2022. \url{https://arxiv.org/abs/2206.06301}
\bibitem{ref131} ENTS: An Edge-native Task Scheduling System for Collaborative Edge Computing. arXiv:2210.07842, 2022. \url{https://arxiv.org/abs/2210.07842}
\bibitem{ref132} Computing Resource Allocation of Mobile Edge Computing Networks Based on Potential Game Theory. arXiv:1901.00233, 2019. \url{https://arxiv.org/abs/1901.00233}
\bibitem{ref133} Energy Efficiency and Delay Tradeoff in an MEC-Enabled Mobile IoT Network. arXiv:2202.03648, 2022. \url{https://arxiv.org/abs/2202.03648}
\bibitem{ref134} KubeDSM: A Kubernetes-based Dynamic Scheduling and Migration Framework for Cloud-Assisted Edge Clusters. arXiv:2501.07130, 2025. \url{https://arxiv.org/abs/2501.07130}
\bibitem{ref135} Heterogeneous FPGA+GPU Embedded Systems: Challenges and Opportunities. arXiv:1901.06331, 2019. \url{https://arxiv.org/abs/1901.06331}
\bibitem{ref136} EdgeAISim: A Toolkit for Simulation and Modelling of AI Models in Edge Computing Environments. arXiv:2310.05605, 2023. \url{https://arxiv.org/abs/2310.05605}
\bibitem{ref137} Optimizing Resource Allocation and Energy Efficiency in Federated Fog Computing for IoT. arXiv:2504.00791, 2025. \url{https://arxiv.org/abs/2504.00791}
\bibitem{ref138} Cooperative Task Offloading through Asynchronous Deep Reinforcement Learning in Mobile Edge Computing for Future Networks. arXiv:2504.17526, 2025. \url{https://arxiv.org/abs/2504.17526}
\bibitem{ref139} EdgeMORE: Improving Resource Allocation with Multiple Options from Tenants. arXiv:1908.01526, 2019. \url{https://arxiv.org/abs/1908.01526}
\bibitem{ref140} Energy Efficient Deployment and Orchestration of Computing Resources at the Network Edge: a Survey on Algorithms, Trends and Open Challenges. arXiv:2209.14141, 2022. \url{https://arxiv.org/abs/2209.14141}
\bibitem{ref141} EdgeMLBalancer: A Self-Adaptive Approach for Dynamic Model Switching on Resource-Constrained Edge Devices. arXiv:2502.06493, 2025. \url{https://arxiv.org/abs/2502.06493}
\bibitem{ref142} Cooling-Aware Resource Allocation and Load Management for Mobile Edge Computing Systems. arXiv:2006.10978, 2020. \url{https://arxiv.org/abs/2006.10978}
\bibitem{ref143} Oakestra white paper: An Orchestrator for Edge Computing. arXiv:2207.01577, 2022. \url{https://arxiv.org/abs/2207.01577}
\bibitem{ref144} Context-aware Container Orchestration in Serverless Edge Computing. arXiv:2408.07536, 2024. \url{https://arxiv.org/abs/2408.07536}
\bibitem{ref145} Online Adaptive Learning for Runtime Resource Management of Heterogeneous SoCs. arXiv:2008.09728, 2020. \url{https://arxiv.org/abs/2008.09728}
\bibitem{ref146} Multiple Resource Allocation in Multi-Tenant Edge Computing via Sub-modular Optimization. arXiv:2302.09888, 2023. \url{https://arxiv.org/abs/2302.09888}
\bibitem{ref147} A Review of Resource Management in Fog Computing: Machine Learning Perspective. arXiv:2209.03066, 2022. \url{https://arxiv.org/abs/2209.03066}
\bibitem{ref148} Local Ratio based Real-time Job Offloading and Resource Allocation in Mobile Edge Computing. arXiv:2503.16794, 2025. \url{https://arxiv.org/abs/2503.16794}
\bibitem{ref149} Edge Intelligence for Energy-efficient Computation Offloading and Resource Allocation in 5G Beyond. arXiv:2011.08442, 2020. \url{https://arxiv.org/abs/2011.08442}
\bibitem{ref150} Redundancy Management for Fast Service (Rates) in Edge Computing Systems. arXiv:2303.00486, 2023. \url{https://arxiv.org/abs/2303.00486}
\bibitem{ref151} Control Aspects for Using RIS in Latency-Constrained Mobile Edge Computing. arXiv:2312.12025, 2023. \url{https://arxiv.org/abs/2312.12025}
\bibitem{ref152} QoS-aware Scheduling of Periodic Real-time Task Graphs on Heterogeneous Pre-occupied MECs. arXiv:2506.12415, 2025. \url{https://arxiv.org/abs/2506.12415}
\bibitem{ref153} Energy-Efficient Federated Edge Learning with Joint Communication and Computation Design. arXiv:2003.00199, 2020. \url{https://arxiv.org/abs/2003.00199}
\bibitem{ref154} Computation Rate Maximization for Wireless Powered Edge Computing With Multi-User Cooperation. arXiv:2402.16866, 2024. \url{https://arxiv.org/abs/2402.16866}
\bibitem{ref155} Resource Allocation for Semantic-Aware Mobile Edge Computing Systems. arXiv:2309.11736, 2023. \url{https://arxiv.org/abs/2309.11736}
\bibitem{ref156} Energy-Efficient Joint Offloading and Wireless Resource Allocation Strategy in Multi-MEC Server Systems. arXiv:1803.07243, 2018. \url{https://arxiv.org/abs/1803.07243}
\bibitem{ref157} Joint Optimization of Signal Design and Resource Allocation in Wireless D2D Edge Computing. arXiv:2002.11850, 2020. \url{https://arxiv.org/abs/2002.11850}
\bibitem{ref158} AI-oriented Medical Workload Allocation for Hierarchical Cloud/Edge/Device Computing. arXiv:2002.03493, 2020. \url{https://arxiv.org/abs/2002.03493}
\bibitem{ref159} Industrial Edge-based Cyber-Physical Systems -- Application Needs and Concerns for Realization. arXiv:2201.00242, 2022. \url{https://arxiv.org/abs/2201.00242}
\bibitem{ref160} Spatio-temporal Bayesian Learning for Mobile Edge Computing Resource Planning in Smart Cities. arXiv:2103.07814, 2021. \url{https://arxiv.org/abs/2103.07814}
\bibitem{ref161} Edge-centric Optimization of Multi-modal ML-driven eHealth Applications. arXiv:2208.02597, 2022. \url{https://arxiv.org/abs/2208.02597}
\bibitem{ref162} Container Orchestration in Edge and Fog Computing Environments for Real-Time IoT Applications. arXiv:2203.05161, 2022. \url{https://arxiv.org/abs/2203.05161}
\bibitem{ref163} Edge AI: On-Demand Accelerating Deep Neural Network Inference via Edge Computing. arXiv:1910.05316, 2019. \url{https://arxiv.org/abs/1910.05316}
\bibitem{ref164} Large Language Models Empowered Autonomous Edge AI for Connected Intelligence. arXiv:2307.02779, 2023. \url{https://arxiv.org/abs/2307.02779}
\bibitem{ref165} Evolution of Convolutional Neural Network (CNN): Compute vs Memory bandwidth for Edge AI. arXiv:2311.12816, 2023. \url{https://arxiv.org/abs/2311.12816}
\bibitem{ref166} Green Edge AI: A Contemporary Survey. arXiv:2312.00333, 2023. \url{https://arxiv.org/abs/2312.00333}
\bibitem{ref167} Prioritized Information Bottleneck Theoretic Framework with Distributed Online Learning for Edge Video Analytics. arXiv:2409.00146, 2024. \url{https://arxiv.org/abs/2409.00146}
\bibitem{ref168} Edge Intelligence: The Confluence of Edge Computing and Artificial Intelligence. arXiv:1909.00560, 2019. \url{https://arxiv.org/abs/1909.00560}
\bibitem{ref169} AI Tax: The Hidden Cost of AI Data Center Applications. arXiv:2007.10571, 2020. \url{https://arxiv.org/abs/2007.10571}
\bibitem{ref170} Edge AI for Internet of Energy: Challenges and Perspectives. arXiv:2311.16851, 2023. \url{https://arxiv.org/abs/2311.16851}
\bibitem{ref171} Onboard Optimization and Learning: A Survey. arXiv:2505.08793, 2025. \url{https://arxiv.org/abs/2505.08793}
\bibitem{ref172} Edge Computing Performance Amplification. arXiv:2305.16175, 2023. \url{https://arxiv.org/abs/2305.16175}
\bibitem{ref173} Scheduling Real-time Deep Learning Services as Imprecise Computations. arXiv:2011.01112, 2020. \url{https://arxiv.org/abs/2011.01112}
\bibitem{ref174} Resource Provisioning in Edge Computing for Latency Sensitive Applications. arXiv:2201.11837, 2022. \url{https://arxiv.org/abs/2201.11837}
\bibitem{ref175} DeepFT: Fault-Tolerant Edge Computing using a Self-Supervised Deep Surrogate Model. arXiv:2212.01302, 2022. \url{https://arxiv.org/abs/2212.01302}
\bibitem{ref176} Federated Learning: Strategies for Improving Communication Efficiency. arXiv:1610.05492, 2016. \url{https://arxiv.org/abs/1610.05492}
\bibitem{ref177} FedLPS: Heterogeneous Federated Learning for Multiple Tasks with Local Parameter Sharing. arXiv:2402.08578, 2024. \url{https://arxiv.org/abs/2402.08578}
\bibitem{ref178} FedHe: Heterogeneous Models and Communication-Efficient Federated Learning. arXiv:2110.09910, 2021. \url{https://arxiv.org/abs/2110.09910}
\bibitem{ref179} On-device Federated Learning with Flower. arXiv:2104.03042, 2021. \url{https://arxiv.org/abs/2104.03042}
\bibitem{ref180} EdgeML: Towards Network-Accelerated Federated Learning over Wireless Edge. arXiv:2111.09410, 2021. \url{https://arxiv.org/abs/2111.09410}
\bibitem{ref181} FedCache: A Knowledge Cache-driven Federated Learning Architecture for Personalized Edge Intelligence. arXiv:2308.07816, 2023. \url{https://arxiv.org/abs/2308.07816}
\bibitem{ref182} FedGreen: Federated Learning with Fine-Grained Gradient Compression for Green Mobile Edge Computing. arXiv:2111.06146, 2021. \url{https://arxiv.org/abs/2111.06146}
\bibitem{ref183} High-Dimensional Stochastic Gradient Quantization for Communication-Efficient Edge Learning. arXiv:1910.03865, 2019. \url{https://arxiv.org/abs/1910.03865}
\bibitem{ref184} To Talk or to Work: Flexible Communication Compression for Energy Efficient Federated Learning over Heterogeneous Mobile Edge Devices. arXiv:2012.11804, 2020. \url{https://arxiv.org/abs/2012.11804}
\bibitem{ref185} Communication-Computation Trade-Off in Resource-Constrained Edge Inference. arXiv:2006.02166, 2020. \url{https://arxiv.org/abs/2006.02166}
\bibitem{ref186} FedIN: Federated Intermediate Layers Learning for Model Heterogeneity. arXiv:2304.00759, 2023. \url{https://arxiv.org/abs/2304.00759}
\bibitem{ref187} Scalable and Communication-efficient Decentralized Federated Edge Learning with Multi-blockchain Framework. arXiv:2008.04743, 2020. \url{https://arxiv.org/abs/2008.04743}
\bibitem{ref188} Agglomerative Federated Learning: Empowering Larger Model Training via End-Edge-Cloud Collaboration. arXiv:2312.11489, 2023. \url{https://arxiv.org/abs/2312.11489}
\bibitem{ref189} EdgeFL: A Lightweight Decentralized Federated Learning Framework. arXiv:2309.02936, 2023. \url{https://arxiv.org/abs/2309.02936}
\bibitem{ref190} Providing Location Information at Edge Networks: A Federated Learning-Based Approach. arXiv:2205.08292, 2022. \url{https://arxiv.org/abs/2205.08292}
\bibitem{ref191} Split Federated Learning Empowered Vehicular Edge Intelligence: Concept, Adaptive Design and Future Directions. arXiv:2406.15804, 2024. \url{https://arxiv.org/abs/2406.15804}
\bibitem{ref192} Uncertainty Estimation in Multi-Agent Distributed Learning for AI-Enabled Edge Devices. arXiv:2403.09141, 2024. \url{https://arxiv.org/abs/2403.09141}
\bibitem{ref193} Cognitive Edge Computing: A Comprehensive Survey on Optimizing Large Models and AI Agents for Pervasive Deployment. arXiv:2501.03265, 2025. \url{https://arxiv.org/abs/2501.03265}
\bibitem{ref194} Integrated Sensing-Communication-Computation for Edge Artificial Intelligence. arXiv:2306.01162, 2023. \url{https://arxiv.org/abs/2306.01162}
\bibitem{ref195} Roadmap for Edge AI: A Dagstuhl Perspective. arXiv:2112.00616, 2021. \url{https://arxiv.org/abs/2112.00616}
\bibitem{ref196} Enabling All In-Edge Deep Learning: A Literature Review. arXiv:2204.03326, 2022. \url{https://arxiv.org/abs/2204.03326}
\bibitem{ref197} Edge Intelligence for Empowering IoT-Based Healthcare Systems. arXiv:2103.12144, 2021. \url{https://arxiv.org/abs/2103.12144}
\bibitem{ref198} Convergence of Edge Computing and Deep Learning: A Comprehensive Survey. arXiv:1907.08349, 2019. \url{https://arxiv.org/abs/1907.08349}
\bibitem{ref199} On the Sustainability of AI Inferences in the Edge. arXiv:2507.23093, 2025. \url{https://arxiv.org/abs/2507.23093}
\bibitem{ref200} DAG-based Task Orchestration for Edge Computing. arXiv:2301.09278, 2023. \url{https://arxiv.org/abs/2301.09278}
\bibitem{ref201} Resilient Edge: Can we achieve Network Resiliency at the IoT Edge using LPWAN and WiFi?. arXiv:2205.03729, 2022. \url{https://arxiv.org/abs/2205.03729}
\bibitem{ref202} EdgeToll: A Blockchain-based Toll Collection System for Public Sharing of Heterogeneous Edges. arXiv:1912.12681, 2019. \url{https://arxiv.org/abs/1912.12681}
\bibitem{ref203} Sustainable Edge Computing: Challenges and Future Directions. arXiv:2304.04450, 2023. \url{https://arxiv.org/abs/2304.04450}
\bibitem{ref204} Mobility Aware Edge Computing Segmentation Towards Localized Orchestration. arXiv:2110.07808, 2021. \url{https://arxiv.org/abs/2110.07808}
\bibitem{ref205} Toward Edge-enabled Cyber-Physical Systems Testbeds. arXiv:1910.01173, 2019. \url{https://arxiv.org/abs/1910.01173}
\bibitem{ref206} Network Emulation in Large-Scale Virtual Edge Testbeds: A Note of Caution and the Way Forward. arXiv:2208.05862, 2022. \url{https://arxiv.org/abs/2208.05862}
\bibitem{ref207} 6G White Paper on Edge Intelligence. arXiv:2004.14850, 2020. \url{https://arxiv.org/abs/2004.14850}
\bibitem{ref208} 6G-enabled Edge AI for Metaverse: Challenges, Methods, and Future Research Directions. arXiv:2204.06192, 2022. \url{https://arxiv.org/abs/2204.06192}
\bibitem{ref209} Mobile Edge Computing, Metaverse, 6G Wireless Communications, Artificial Intelligence, and Blockchain: Survey and Their Convergence. arXiv:2209.14147, 2022. \url{https://arxiv.org/abs/2209.14147}
\bibitem{ref210} EdgeSphere: A Three-Tier Architecture for Cognitive Edge Computing. arXiv:2405.16685, 2024. \url{https://arxiv.org/abs/2405.16685}
\bibitem{ref211} SD-VEC: Software-Defined Vehicular Edge Computing with Ultra-Low Latency. arXiv:2103.14225, 2021. \url{https://arxiv.org/abs/2103.14225}
\bibitem{ref212} MEET: Mobility-Enhanced Edge inTelligence for Smart and Green 6G Networks. arXiv:2210.15111, 2022. \url{https://arxiv.org/abs/2210.15111}
\bibitem{ref213} Key Enabling Technologies for Secure and Scalable Future Fog-IoT Architecture: A Survey. arXiv:1806.06188, 2018. \url{https://arxiv.org/abs/1806.06188}
\bibitem{ref214} A Survey on Edge Computing Systems and Tools. arXiv:1911.02794, 2019. \url{https://arxiv.org/abs/1911.02794}
\bibitem{ref215} Quality of Service (QoS)-driven Edge Computing and Smart Hospitals: A Vision, Architectural Elements, and Future Directions. arXiv:2303.06896, 2023. \url{https://arxiv.org/abs/2303.06896}
\bibitem{ref216} Ubiquitous Computing: Potentials and Challenges. arXiv:1011.1960, 2010. \url{https://arxiv.org/abs/1011.1960}
\bibitem{ref217} A Quality-of-Service Compliance System using Federated Learning and Optimistic Rollups. arXiv:2312.00026, 2023. \url{https://arxiv.org/abs/2312.00026}
\bibitem{ref218} Challenges and Opportunities in Edge Computing. arXiv:1609.01967, 2016. \url{https://arxiv.org/abs/1609.01967}
\bibitem{ref219} Future Computer Systems and Networking Research in the Netherlands: A Manifesto. arXiv:2206.03259, 2022. \url{https://arxiv.org/abs/2206.03259}
\end{thebibliography}

\end{document}
